\documentclass[11pt,a4paper]{article}
\usepackage[margin=1.8cm]{geometry}
\usepackage{graphicx}\usepackage{amsmath,amssymb,tabularx,array,needspace,fontspec,enumitem,cancel}
\newcommand{\tb}{\cancel{||||}\,}
\setmainfont{FreeSerif}
\newfontfamily\urdufont[Script=Arabic]{FreeSerif}
\TeXXeTstate=1
\newcommand{\urdu}[1]{{\urdufont\beginR #1\endR}}
\hyphenpenalty=10000\exhyphenpenalty=10000
\setlength{\parindent}{0pt}
\setlength{\parskip}{0pt}
\renewcommand{\arraystretch}{1.15}
\pagestyle{plain}

% \q{roman}{text}
\newlength{\qw}\settowidth{\qw}{\textbf{888.}\ }
\newcommand{\q}[2]{%
  \Needspace{3\baselineskip}%
  \par\medskip\noindent\hangindent=\qw\hangafter=1\makebox[\qw][l]{\textbf{#1.}}#2\par}

\begin{document}

\begin{center}
  {\Huge\bfseries The Hope Science Academy}\\[3pt]
  {\small Excellence is our standard not goal}\\[2pt]
  {\small 247-E-1, (Opposite LDA School) Johar Town, Lahore. PH: 0300-8194789, 0322-4157001}
\end{center}

\medskip
\noindent
\begin{tabularx}{\linewidth}{|>{\raggedright\arraybackslash}X|>{\raggedright\arraybackslash}X|>{\raggedright\arraybackslash}X|>{\raggedright\arraybackslash}X|}
\hline
\textbf{Student Name:} & \textbf{Roll Num:} & \textbf{Class:} 12th & \textbf{Paper Code:} \\ \hline
\textbf{Subject Name:} Mathematics. & \textbf{Time Allowed:} & \textbf{Total Marks:} 977 & \textbf{Exam Date:} 03-Nov-2026 \\ \hline
\multicolumn{4}{|l|}{\textbf{Exam Syllabus:} FULL BOOK} \\ \hline
\end{tabularx}

\bigskip
\noindent\textbf{Q2. Write detailed answers of the following questions.}\hfill\textbf{1$\times$603=603}

\q{1}{Discuss domain, range and graph of the function $f(x)=3^{x}$}
\q{2}{Discuss domain, range and graph of the function $f(x)=\left(\dfrac{1}{3}\right)^{x}$}
\q{3}{Prove the following inverse hyperbolic formulae.\par $\tanh^{-1}x=\dfrac{1}{2}\ln\!\left(\dfrac{1+x}{1-x}\right),x\neq1$}
\q{4}{Prove the following inverse hyperbolic formulae.\par $\coth^{-1}x=\dfrac{1}{2}\ln\!\left(\dfrac{1+x}{1-x}\right),x\neq1$}
\q{5}{Prove the following inverse hyperbolic formulae.\par $\mathrm{sech}^{-1}x=\dfrac{1}{2}\ln\!\left(\dfrac{1+\sqrt{1-x^{2}}}{x}\right),x\neq1$}
\q{6}{Prove the following inverse hyperbolic formulae.\par $\mathrm{csch}^{-1}x=\ln\!\left(\dfrac{1}{x}+\sqrt{\dfrac{1}{x^{2}}+1}\right)$}
\q{7}{Given the following function classify each function as algebraic or transcendental.\par $f(x)=\sqrt{x+1},\,g(x)=e^{x},\,h(x)=\dfrac{x+1}{x^{2}-2x+1},\,q(x)=\ln 2x,\ r(x)=\sin x$}
\q{8}{Graph of $f(x)=e^{x}$}
\q{9}{Graph of $f(x)=\left(\dfrac{1}{2}\right)^{x}=\dfrac{1}{2^{x}}$}
\q{10}{Prove that: $\sinh^{-1}x=\ln\!\left(x+\sqrt{x^{2}+1}\right)$}
\q{11}{Prove that: $\cosh^{-1}x=\ln\!\left(x+\sqrt{x^{2}-1}\right)$ for all $x\in[1,\infty)$}
\q{12}{Let $f(x)=\log_{4}(x-2)$ : show that f(x) is a one-to-one function using the horizonta line test.}
\q{13}{Let $f(x)=\log_{4}(x-2)$ : Find the inverse function $f^{-1}(x)$ algebraically.}
\q{14}{Let $f(x)=\log_{4}(x-2)$ : State the domain and range of f(x) and $f^{-1}(x)$.}
\q{15}{Let $f(x)=\log_{4}(x-2)$ : Verify graphically that the graphs of f(x) and Let $f^{-1}(x)$ are symmetric about the line y=x.}
\q{16}{Let $f(x)=\log_{4}(x-2)$ : Verify graphically that the graphs of f(x) and Let $f^{-1}(x)$ are symmetric about the line y=x.}
\q{17}{Let $f(x)=\sin x$ for all $x\in\left[-\dfrac{\pi}{2},\dfrac{\pi}{2}\right]$ : state the domain and range of $f(x)$ and $f^{-1}(x)$}
\q{18}{Let $f(x)=\sin x$ for all $x\in\left[-\dfrac{\pi}{2},\dfrac{\pi}{2}\right]$ : State the graph of f(x) and show that it is one-to-one using the horizontal line test.}
\q{19}{Let $f(x)=\sin x$ for all $x\in\left[-\dfrac{\pi}{2},\dfrac{\pi}{2}\right]$ : Draw the graph of Let $f^{-1}(x)$ using the reflection of f(x) along the line y=x :}
\q{20}{A one-to one function is defined by $f(x)=(x-1)^{3},\ x\in[0,3]$ using the concept of reflection along the line y=x, sketch the graph of the inverse function $f^{-1}(x)$.}
\q{21}{Given the function $f(x)=\sqrt{x}$ find the equation of the function after: Applying both transformations\par (i) vertical stretch by factor 2 (ii) Horizontal shift 3 units to the right.}
\q{22}{Given $f(x)=\dfrac{1}{x},x\neq0$ the function find the equation of the transformed function after applying the following transformation in order. Horizontal shift 2 units to the left, horizontal compression by factor 2, vertical stretching by factor 3, and vertical shift 1 unit upward.}
\q{23}{Given the function f(x)= ln x, find the equation of the transformed function obtained after shifting the graph 1 unit to the right, vertically compressing it by a factor of 2, and then shifting it 3 units upward.}
\q{24}{Given the function $f(x)=\sqrt{\dfrac{2x-1}{x+3}}$ . Find the equation of the transformed function after applying in order, the horizontal shift 2 units to the left, horizontal, compression by factor 3, vertical compression by factor 2, and vertical shift 1 unit downward.}
\q{25}{Given the function $f(x)=3^{x}$ and its transformed function $g(x)=3^{x-2}+1$ , draw the graphs of both functions on the same coordinate plane and compare their shapes.}
\q{26}{Given the function $f(x)=\ln x$ . Find the equation of the transformed function after applying in order, the vertical stretch by factor 4, shift 6 units to right, horizontal compress by factor 3, and shift 1 unit upward. Draw the graphs of both the original and transformed functions on the same coordinate plane.}
\q{27}{Draw the graph of $f(x)=|x|$ after applying in order, shift of 7 units to right horizontal compress by factor 2, vertical shift 6 units downward, and vertical copress by factor 2.}
\q{28}{Show graphically that the function y= ln x is the inverse of the function $y=e^{x}$ .}
\q{29}{Let $f(x)=\dfrac{x}{x+1}$ . find the function obtained after applying all the following transformations with the given order: (i) Verictal stretch by 2 (ii) Horizontal stretch by 3.}
\q{30}{Let $f(x)=\dfrac{x}{x+1}$ . find the function obtained after applying all the following transformations with the given order: (i) Shift 2 units to the right (ii) Shift 1 unit downward}
\q{31}{Solve the following equation: $\log(x)+\log(x^{2}-5x+7)=\log(3)$}
\q{32}{Solve the following equation: $3^{(x+1)}-2.9^{x}+9=0$}
\q{33}{Solve the following inequalitie (if its solution exists): $\log_{\frac{1}{3}}(x^{2}-4x+3)>\log_{\frac{1}{3}}(2x-1)$}
\q{34}{Solve the following inequalitie (if its solution exists): $\log_{2}(x-2)+\log_{2}(x-3)\geq3$}
\q{35}{Solve the following inequalitie (if its solution exists): $\log(x)+\log(x-3)\geq\log(2x)+\log(x-1)$}
\q{36}{Find the solution of $2^{2x}-3^{5x+3}=0$}
\q{37}{Solve the equation: $\log x+\log(x-3)=1$}
\q{38}{Solve the equation: $\log_{x}(x^{2}-3x+2)=2$}
\q{39}{Solve the inequality: $2^{x+1}>5^{x}$}
\q{40}{Solve the inequality: $\log_{\frac{7}{9}}(x^{2}-3x+2)\geq\log_{7}(2x-1)$}
\q{41}{A quantity chnages according to $N=500e^{kt}$ . If N=250 when t=5, find: (i) the value of constant k. (ii) The quantity N when t=10.}
\q{42}{A sum of Rs 5000 becomes 6324.48 when invested at 8\% per annum, compounded annually. Find the time t.}
\q{43}{A sum of Rs 3000 is invested at 5\% per annum, compounded continously for 4 years. Find the final amount.}
\q{44}{An amount of Rs 2000 becomes 3320.12 under continuous compounding at 6\% per annum. Find the time t.}
\q{45}{The sound levels at two points are L$_{1}$ =90dB and L$_{2}$ =60dB. Find the ratio of their intensities $\dfrac{I_{1}}{I_{2}}$, where $I_{o}10^{-12}W/m^{2}$ .}
\q{46}{A sum of Rs 1000 is invested at a rate of 5\% per year compounded continuously for 3 years. Find the final amount.}
\q{47}{A sum of Rs 5000 is invested at a rate of 8\% per annum, compounded quarterly. Find the amount aafter 3 years.}
\q{48}{Derive the formula for the derivative of the following function w.r.t. x: sec x}
\q{49}{Derive the formula for the derivative of the following function w.r.t. x: csc x.}
\q{50}{Derive the formula for the derivative of the following function w.r.t. x: cot x.}
\q{51}{Find $\dfrac{dy}{dx}$ by making suitable substitutions and applying the China Rule.\par $y=\sqrt{\dfrac{1+x}{1-x}}$}
\q{52}{Find $\dfrac{dy}{dx}$ by making suitable substitutions and applying the China Rule.\par $y=(\sin x+\cos x)^{3/2}$}
\q{53}{Find $\dfrac{dy}{dx}$ for the following parametric equation:\par $x=\sqrt{1+\sec t}$ , and $y=\sqrt{\tan t}$}
\q{54}{If $x=a\sec t$ and $y=b\tan t$, then show that $\dfrac{dy}{dx}=\dfrac{b^{2}x}{a^{2}x}$}
\q{55}{Use implicit differentiation to find $\dfrac{dy}{dx}$ if:\par $ax^{2}+by^{2}+2gx+2fy+c=0$}
\q{56}{Differentiate w.r.t the variables involved: $\csc^{-1}\!\left(\dfrac{x^{2}+1}{x^{2}-1}\right)$}
\q{57}{Differentiate w.r.t the variables involved: $\cot^{-1}(\sqrt{x})+\sqrt{\cot^{-1}x}$}
\q{58}{Differentiate w.r.t the variables involved: $(\tan)^{-1}+\dfrac{1}{\tan^{-1}t}$}
\q{59}{Find $\dfrac{dy}{dx}$ if $y=a\cos t$ and $x=a\sin t$.}
\q{60}{Find $\dfrac{dy}{dx}$ if $y=x\cos^{-1}\!\left(\dfrac{x}{a}\right)-\sqrt{a^{2}-x^{2}},\ a>0.$}
\q{61}{Find $\dfrac{dy}{dx}$ if: $y=\log_{10}\sqrt{\dfrac{1-x^{2}}{1+x^{2}}}$}
\q{62}{Find $\dfrac{dy}{dx}$ if: $y=\dfrac{\sqrt{x^{2}-1}\,(2x+1)}{(x^{3}-1)^{3/2}}$}
\q{63}{Find the equation of the tangent line to the curve $y=\sin^{-1}\!\left(\dfrac{x}{2}\right)+\cos x$ at the point where $x=\dfrac{\pi}{2}$.}
\q{64}{Find a point on the curve $y=\tan^{-1}\!\left(\dfrac{x}{3}\right)-\dfrac{x}{3}$ , where tangent is horizontal. Also find the equation of the tangent and normal at that point.}
\q{65}{Find a point on the curve $x+y-e^{y}=0$ where the tangent line is vertical. Also find the equations of the tangent and normal at that point.}
\q{66}{Find all points on the ellipse $x^{2}-xy+y^{2}=1$ at which the tangent line is horizontal.}
\q{67}{Find $y_{2}$ if: $x=a\cos t,\ y=b\sin t$}
\q{68}{Find $y_{2}$ if: $x=t^{2},\ y=2t^{3}$}
\q{69}{ind $y_{4}$ and write expressions for $y_{10}$ and $y_{11}$ without evaluating them: $y=\sin\!\left(\dfrac{x}{2}\right)$}
\q{70}{Find $y_{4}$ and write expressions for $y_{10}$ and $y_{11}$ without evaluating them: $y=\ln(x^{2}-a^{2})$}
\q{71}{If $x^{2}+y^{2}=25$, then show that $y^{3}\dfrac{d^{2}y}{dx^{2}}+25=0.$}
\q{72}{If $y=e^{ax}\cos bx$ , then show that $\dfrac{d^{2}y}{dx^{2}}-2a\dfrac{dy}{dx}+(a^{2}+b^{2})y=0.$}
\q{73}{If $x=\sin t$ and $y=\cos^{2}t$, then show that $\dfrac{d^{2}y}{dx^{2}}+2=0.$}
\q{74}{If x= ln t and $y=\dfrac{1}{t}$, then show that $\dfrac{d^{2}y}{dx^{2}}-y=0.$}
\q{75}{If $y=a\cos(\ln x)+b\sin(\ln x)$, then show that $x^{2}y_{2}+xy_{1}+y=0.$}
\q{76}{Find the equations of the tangent and normal to the curve y=sec x at the point where $x=\dfrac{\pi}{4}$ .}
\q{77}{Find the higher- order derivatives of the function\par $f(x)=\dfrac{1}{4}x^{4}-x^{3}+\dfrac{1}{2}x^{2}-3$}
\q{78}{Find all the higher-order derivatives of the function $y=e^{ax}$ .}
\q{79}{If $x=e^{at}$ and $y=e^{2at}$ show that $x^{2}\dfrac{d^{2}y}{dx^{2}}-2y=0.$}
\q{80}{Find $y_{4}$ if $y=\ln\sqrt{\dfrac{2x-3}{x+1}}$ .}
\q{81}{Determine the intervals on which f is increasing or decreasing for the domain mentioned in following: $f(x)=\sin x\ ;\ x\in(0,2\pi)$}
\q{82}{Determine the intervals on which f is increasing or decreasing for the domain mentioned in following: $f(x)=4-x^{2}\ ;\ x\in R=(-\infty,\infty)$}
\q{83}{Determine the intervals on which f is increasing or decreasing for the domain mentioned in following: $f(x)=x^{2}+3x+2\ ;\ x\in[-4,1]$}
\q{84}{Determine the intervals on which f is increasing or decreasing for the domain mentioned in following: $f(x)=\sin x+\dfrac{1}{4}\cos 2x\ ;\ x\in[0,2\pi]$}
\q{85}{Using the second derivative test, find local extreme values of the following function: $f(x)=8-2x^{2}$}
\q{86}{Using the second derivative test, find local extreme values of the following function: $f(x)=x^{2}-x-2$}
\q{87}{Using the second derivative test, find local extreme values of the following function: $f(x)=x^{4}-8x^{2}+2$}
\q{88}{Using the second derivative test, find local extreme values of the following function: $f(x)=(x+2)^{2}(x-3)$}
\q{89}{Using the second derivative test, find local extreme values of the following function: $f(x)=(x^{2}-2x)e^{x}$}
\q{90}{Discuss the global extrema of the following function over specified domains: $f(x)=\sin x+\cos x$ on $[0,2\pi]$}
\q{91}{Discuss the global extrema of the following function over specified domains: $f(x)=x^{4/3}$ on $[-1,8]$}
\q{92}{Find the global and local extrema of the following functions over thr specified domains (hint) use the increasing and decrasing behavior of the functions including the end points). $f(x)=4-x^{2}$ on $[-1,2]$}
\q{93}{Find the global and local extrema of the following functions over thr specified domains (hint) use the increasing and decrasing behavior of the functions including the end points). $f(x)=\cos x-\sin x$ on $[0,2\pi]$}
\q{94}{Prove that $f(x)=\dfrac{\ln(2x)}{x^{2}}$ has local maximum at $x=\dfrac{\sqrt{e}}{2}$ .}
\q{95}{Find the intervals on which the function $f(x)=|x|$ , defined on $R=(-\infty,\infty)$, is increasing or decreasing.}
\q{96}{Find the intervals in which the function $f(x)=2x^{3}-3x^{2}-12x$ is increasing or decreasing.}
\q{97}{Find the intervals where the function $f(x)=x^{\frac{2}{3}}$ is increasing or decreasing on the domain $R=(-\infty,\infty)$}
\q{98}{Find the global extreme values of $f(x)=x^{4}-4x$ on $[-2,3]$.}
\q{99}{Examine the function $f(x)=\dfrac{1}{2}x^{4}-\dfrac{4}{3}x^{3}-x^{2}+4x+1$ for relative extreme values.}
\q{100}{Find the linear approximation of $f(x)=\sqrt{x+1}$ at $x=3$ and use it to approximate $\sqrt{3{,}1}$ and $\sqrt{2{,}9}$ .}
\q{101}{Find the local linear approximation of $f(x)=\sin x$ at $x=0$.}
\q{102}{Use the local linear approximation obtained in $f(x)=\sin x$ at $x=0$ to approximate sin $\sin 2^{\circ}$ , and compare your approximation to the result produced directed by a calculator.}
\q{103}{Find the linear approximation of the following function at the given points: $f(x)=\sqrt{x^{2}+16}$ at $x=3$}
\q{104}{Find the linear approximation of the following function at the given points: $f(x)=\tan x$ at $x=\pi$}
\q{105}{Find the linear approximation of the following function at the given points: $f(x)=\cos x$ at $x=0$}
\q{106}{Use a local linear approximation to estimate the following: $(82)^{\frac{1}{3}}$}
\q{107}{Use a local linear approximation to estimate the following: $\sin 29^{\circ}$}
\q{108}{Find the local linear approximation of $f(x)=\sqrt{2x-1}$ at $x=5$.}
\q{109}{Use the local linear approximation obtained in $f(x)=\sqrt{2x-1}$ at $x=5$ to approximate $\sqrt{9.2}$ and compare your approximation to the result produced directly by calculator.}
\q{110}{Find the local linear approximation of $f(x)=\cos x$ at $x=\dfrac{\pi}{2}$ .}
\q{111}{Use the local linear approximatiob obtained in $f(x)=\cos x$ at $x=\dfrac{\pi}{2}$ to approximate cos 91 degree, and compare your approximation to the result produced directly by a calculator.}
\q{112}{Use dy to approximate $\triangle y$ for the following function: $y=\sin x+\cos x$ when x changes from $30^{\circ}$ to $29^{\circ}$.}
\q{113}{Use dy to approximate $\triangle y$ for the following function: $y=\dfrac{1}{x}$ when $x=0.5$ and $dx=-0.1$.}
\q{114}{Use differentials to approximate the value of: $\sqrt{26}$}
\q{115}{Use differentials to approximate the value of: $(33)^{\frac{1}{5}}$}
\q{116}{Use differentials to approximate the value of: $\tan 46^{\circ}$}
\q{117}{The side of a cube is measured to be 30 cm with a masimum error of 0.14 cm in its measurement. Find the maximum error in the calculated volume of the cube.}
\q{118}{The radius of a sphere is measured as $100\pm0.5$ cm and the volume is calculated from this measurement. Estimate the percentage error in the volume calculation.}
\q{119}{Use differentials to approximate $\sqrt{24}$ .}
\q{120}{Using differentials approximate the value of cos $61^{\circ}$ .}
\q{121}{The side of a cube is measured to 25 cm, with a possible error of $\pm1$ cm. Use differentials to estimate the error in the calculated volume.}
\q{122}{The side of a cube is measured to 25 cm, with a possible error of $\pm1$ cm. Estimate the percentage error in the side and volume.}
\q{123}{The radius r of a circle is measured with an error of at most 2\%. What is the maximum corresponding percentage error in computing the circle's.}
\q{124}{A window consist of a rectangle surmounted by a semicircular opening. If the total perimeter is 21 meter, find the dimensions that maximize the area.}
\q{125}{A rectangular sheet of meta with dimensions 40 cm by 25 cm is to be make into an open box by cutting equal squares from the corners and folding up the sides. Find the size of the tank squares that maximize the volume.}
\q{126}{A closed cylindrixal tank is to constructed to hold $54\pi$ liters of water. Find the dimensions that minimize the construction cost of the tank.}
\q{127}{A wire of 20 metres is cut into two pices. One piece is bent into a square and the other into a circle. How should the wire be cut to minimize the total area enclosed?}
\q{128}{A team of environmental  scientists is monitoring the oxygen level in a lake over time to esnure the water quality remains suitable for aquatic life. The oxygen level in the lake, measured in milligram per liter (mg/L) is modeled by the equation, $O(t)=t^{3}-18t^{2}+81t+100$ where O(t) is the oxygn level at time t (in days aftr monitoring began). scientists want to determine when during the monitoring period the oxygen level is at its lowest, as low oxygen levels can be harmful to fish and other quatic organisms. Find the time t when the oxygen level is at its minimum.}
\q{129}{The perimeter of a triangle is 16 cm. If one side of length 6 cm, what are the length of the other sides for maximum area of the triangle.}
\q{130}{A poster must have 200 cm2 printed area with margins of 2 cm on top/ bottom and 3 cm on sides. Find poster dimensions to minimize papaer used.}
\q{131}{Find two positive integers whose sum is 12 and the product of one with the square of the other will be maximum.}
\q{132}{Find the dimensions of a rectange of maximum area that can be incribed under the curve $y=4-x^{2}$ and above the x-axis.}
\q{133}{A cylindrical can must hold $500\ cm^{3}$ of liquid. Find the radius and height that minimizes the material cost.}
\q{134}{A car moves along the parabolic road $y=4-x^{2}$. A mobile tower is at the origin. Find the position of the car where the signal strength is strongest.}
\q{135}{A software company is designing a data processing algorithm or analyzing batches of customer transactions. The runtime T(n) in milliseconds depends on the size of the batch n (the umber of transaction processed together) and is modeled by the equation.}
\q{136}{Evaluate the following: $\displaystyle\int\dfrac{x+1}{\sqrt{x+2}}\,dx$}
\q{137}{Evaluate the following: $\displaystyle\int\dfrac{1}{\sqrt{x+1}+\sqrt{x}}\,dx$}
\q{138}{Evaluate the following: $\displaystyle\int\dfrac{a^{x}+a^{-x}}{a^{x}-a^{-x}}\,dx\ \left(\begin{matrix}a>0\\a\neq1\end{matrix}\right)$}
\q{139}{Evaluate the following: $\displaystyle\int\dfrac{2^{2x}+1}{x^{x}}\,dx$}
\q{140}{Evaluate the following: $\displaystyle\int\dfrac{x^{2}+1}{x^{3}-x}\,dx$}
\q{141}{Evaluate the following: $\displaystyle\int\sec^{4}x\,dx$}
\q{142}{Evaluate the following: $\displaystyle\int\dfrac{\cos^{3}x-\sin^{3}x}{\cos x-\sin x}\,dx$}
\q{143}{Evaluate the following: $\displaystyle\int\dfrac{\cos 2x}{\sin x+\cos x}\,dx$}
\q{144}{Evaluate the following: $\displaystyle\int\dfrac{1}{1+\cos x}\,dx$}
\q{145}{Evaluate the following: $\displaystyle\int\dfrac{1-\cos 2x}{1+\cos 2x}\,dx$}
\q{146}{Evaluate the following: $\displaystyle\int\dfrac{\cos 2x-\cos 2\theta}{\cos x-\cos\theta}\,dx$}
\q{147}{Evaluate the integrals: $\displaystyle\int\dfrac{(1+x^{2})^{2}}{x^{4}}\,dx$}
\q{148}{Evaluate: $\displaystyle\int\dfrac{\tan x-\sec x}{1+\tan^{2}x}\,dx$}
\q{149}{Evaluate: $\displaystyle\int\dfrac{\sin(x+a)}{\sin(x+b)}\,dx$}
\q{150}{Evaluate $\displaystyle\int\dfrac{a^{2x}-b^{2x}}{(ab)^{x}}$ dx where a and b are distinct and have the same sign.}
\q{151}{Evaluate the following integrals: $\displaystyle\int x^{5}\sqrt{x^{3}+1}\,dx$}
\q{152}{Evaluate the following integrals: $\displaystyle\int\dfrac{1}{1+\sqrt{2+x}}\,dx$}
\q{153}{Evaluate the following integrals: $\displaystyle\int\dfrac{\sin x-\cos x}{\sqrt{1+\sin 2x}}$}
\q{154}{Evaluate the following integrals: $\displaystyle\int\dfrac{\sin 2x}{\sqrt{1+\sin^{2}x}}\,dx$}
\q{155}{Evaluate the following integral: $\displaystyle\int\dfrac{\cos x}{\sin^{2}x-4}\,dx$}
\q{156}{Evaluate the following integral: $\displaystyle\int\dfrac{x^{2}}{x^{2}+a^{2}}\,dx$}
\q{157}{Evaluate the following integral: $\displaystyle\int\dfrac{1}{x^{2}\sqrt{x^{2}+4}}\,dx$}
\q{158}{Evaluate the following integral: $\displaystyle\int\dfrac{1}{(1+x^{2})^{\frac{3}{2}}}\,dx$}
\q{159}{Evaluate the following integral: $\displaystyle\int\dfrac{x^{2}+1}{x\sqrt{x^{2}-9}}\,dx$}
\q{160}{Evaluate the following integral: $\displaystyle\int\sqrt{\dfrac{1+x}{1-x}}\,dx$}
\q{161}{Evaluate the following integral: $\displaystyle\int\dfrac{\sqrt{a^{2}-x^{2}}}{x}\,dx\,(a>0)$}
\q{162}{Evaluate the following integral: $\displaystyle\int\dfrac{\tan x}{\sec^{4}x+\tan^{4}x}\,dx$}
\q{163}{Evaluate the following integral: $\displaystyle\int\dfrac{1}{\sqrt{x^{2}+2x+10}}\,dx$}
\q{164}{Integrate: $\displaystyle\int\sqrt{x}\sqrt{1+\sqrt{x}}\;dx$}
\q{165}{Evaluate: $\displaystyle\int\dfrac{x}{\sqrt{1+x^{2}}+1+x^{2}}\,dx$}
\q{166}{Evaluate: $\displaystyle\int\dfrac{e^{2x}+2e^{x}}{(e^{x}+1)^{2}}\,dx$}
\q{167}{Evaluate: $\displaystyle\int\dfrac{\ln(1+x)-\ln x}{x(1+x)}\,dx$}
\q{168}{Evaluate: $\displaystyle\int\dfrac{x^{2}\tan^{-1}(x^{3})}{1+x^{6}}\,dx$}
\q{169}{Evaluate: $\displaystyle\int\dfrac{\sqrt{9-x^{2}}}{x^{2}}\,dx$}
\q{170}{Evaluate: $\displaystyle\int\dfrac{1}{x\sqrt{x^{4}+1}}\,dx,\ x\in(-\infty,0)\cup(0,\infty)$}
\q{171}{Evaluate: $\displaystyle\int\dfrac{\sqrt{x^{2}-9}}{x}\,dx,\ (x>3)$}
\q{172}{Evaluate: $\displaystyle\int\dfrac{1}{x\sqrt{1-x^{3}}}\,dx,\ x\in(-\infty,0)\cup(0,1)$}
\q{173}{Evaluate: $\displaystyle\int\dfrac{dx}{(x-2)\sqrt{x^{2}-4x+3}}$}
\q{174}{Evaluate the following integrals: Evaluate: $\displaystyle\int e^{x^{2}}(2x^{3}+2x)\,dx$}
\q{175}{Evaluate the following integrals: Evaluate: $\displaystyle\int\dfrac{x}{1+\cos x}\,dx$}
\q{176}{Evaluate the following integrals: Evaluate: $\displaystyle\int\dfrac{x}{1-\sin x}\,dx$}
\q{177}{Evaluate the following integrals: Evaluate: $\displaystyle\int\left[2\ln x+(\ln x)^{2}\right]dx$}
\q{178}{Evaluate the following integrals: Evaluate: $\displaystyle\int\dfrac{\tan^{-1}\sqrt{x}}{\sqrt{x}}\,dx$}
\q{179}{Evaluate the following integrals: Evaluate: $\displaystyle\int e^{x}\left(\dfrac{1-x}{1+x^{2}}\right)^{2}dx$}
\q{180}{Evaluate the following integrals: Evaluate: $\displaystyle\int\dfrac{\ln x\,\tan^{-1}(\ln x)}{x}$}
\q{181}{Evaluate the following integral: $\displaystyle\int\dfrac{e^{\tan^{-1}x}}{\sqrt{x}}\,dx$}
\q{182}{Evaluate the following integrals: Evaluate: $\displaystyle\int\csc^{3}x\,dx$}
\q{183}{Evaluate the following integrals: Evaluate: $\displaystyle\int e^{5x}(x^{2}+x^{4}+1)\,dx$}
\q{184}{Evaluate the following integrals: Evaluate: $\displaystyle\int\dfrac{e^{\sqrt{x}}\left(1+\sqrt{x}\right)}{\sqrt{x}}$}
\q{185}{Evaluate the following integrals: Evaluate: $\displaystyle\int e^{3x}(\sin x-3\cos x)\,dx$}
\q{186}{Evaluate $\displaystyle\int\ln(\sin x)\,.\sin 2x\,dx$}
\q{187}{Evaluate $\displaystyle\int\sin^{-1}\!\left(\dfrac{x}{\sqrt{1+x^{2}}}\right)dx$}
\q{188}{Evaluate $\displaystyle\int e^{ax}\sin bx\,dx=\dfrac{1}{\sqrt{a^{2}+b^{2}}}\,e^{ax}\sin\!\left(bx-\tan^{-1}\dfrac{b}{a}\right)+c$}
\q{189}{Evaluate $\displaystyle\int x\sin^{2}x\,dx$}
\q{190}{Evaluate $\displaystyle\int\sqrt{x^{2}-4x+13}\,dx$}
\q{191}{Evaluate $\displaystyle\int\sqrt{8x-x^{2}-7}\,dx,\ |x-4|\leq3$}
\q{192}{Evaluate $\displaystyle\int\sqrt{x^{2}-8x+7}\,dx,|x-4|\geq3$}
\q{193}{Evaluate the following integral: $\displaystyle\int\dfrac{5x+7}{(x+1)(x+2)}\,dx$}
\q{194}{Evaluate the following integral: $\displaystyle\int\dfrac{3x+5}{x^{2}+4x+3}\,dx$}
\q{195}{Evaluate the following integral: $\displaystyle\int\dfrac{x^{3}+2x^{2}-x-1}{x^{2}-1}\,dx$}
\q{196}{Evaluate the following integral: $\displaystyle\int\dfrac{x^{2}-12x-21}{x^{2}-6x-7}\,dx$}
\q{197}{Evaluate the following integral: $\displaystyle\int\dfrac{3-16x}{4+5x-6x^{2}}\,dx$}
\q{198}{Evaluate the following integral: $\displaystyle\int\dfrac{2x^{2}-8x+6}{x(x-1)(x-2)(x-3)}\,dx$}
\q{199}{Evaluate the following integral: $\displaystyle\int\dfrac{x^{2}-3}{(x+2)^{3}}\,dx$}
\q{200}{Evaluate the following integral: $\displaystyle\int\dfrac{e^{x}-1}{e^{x}(e^{x}+1)}\,dx$}
\q{201}{Evaluate the following integral: $\displaystyle\int\dfrac{25(x+2)}{(x-2)(x+3)^{2}}\,dx$}
\q{202}{Evaluate the following integral: $\displaystyle\int\dfrac{x^{3}+4x^{2}+7x+5}{(x+1)^{3}(x+2)}\,dx$}
\q{203}{Evaluate the following integral: $\displaystyle\int\dfrac{x^{2}+2x-1}{(x-1)(x^{2}+1)}\,dx$}
\q{204}{Evaluate the following integral: $\displaystyle\int\dfrac{2x^{3}+x^{2}+2x+4}{(x^{2}+1)(x^{2}+4)}\,dx$}
\q{205}{valuate the following integral: $\displaystyle\int\dfrac{x^{4}+4x^{2}+2x+1}{(x^{2}+1)^{2}(x+1)}\,dx$}
\q{206}{Evaluate the following integral: $\displaystyle\int\dfrac{x^{2}}{(x^{2}-2x+2)^{2}}\,dx$}
\q{207}{Evaluate the following integral: $\displaystyle\int\dfrac{\cos x}{\sin^{3}x-1}\,dx$}
\q{208}{Integreate: $\displaystyle\int\dfrac{3x-2}{x(x-1)(x-2)}\,dx$}
\q{209}{Evaluate $\displaystyle\int\dfrac{6x-10}{(x+1)(x-3)^{2}}\,dx.$}
\q{210}{Find $\displaystyle\int\dfrac{5x^{2}+3x+2}{(x+1)(x^{2}+1)}\,dx.$}
\q{211}{Evaluate $\displaystyle\int\dfrac{x^{3}+2x+1}{(x^{2}+1)^{2}}\,dx$}
\q{212}{Evaluate the following definite integral: $\displaystyle\int_{0}^{1}\dfrac{\ln(1+x)}{1+x^{2}}\,dx$}
\q{213}{Evaluate the following definite integral: $\displaystyle\int_{0}^{2\pi}\dfrac{\sin^{2}x}{1+e^{\cos x}}\,dx$}
\q{214}{Evaluate the following definite integral: $\displaystyle\int_{-1}^{1}f(x)\,dx$, where $f(x)=\begin{cases}5-3x&\text{if }x\leq0\\5+3x&\text{if }x>0\end{cases}$}
\q{215}{Let $f(x)=\begin{cases}x^{2}&\text{if }0\leq x\leq2\\2x+1&\text{if }2<x\leq4\end{cases}$ Here g(x) is defined as: $g(x)=f(x)+f(4-x)$ compute $\displaystyle\int_{0}^{4}g(x)\,dx$}
\q{216}{Simplify: $\displaystyle\int_{0}^{3}5f(x)\,dx-\int_{1}^{3}5f(4-x)\,dx+\int_{1}^{5}5f(x)\,dx-\int_{0}^{5}4f(5-x)\,dx$}
\q{217}{Show that $\displaystyle\int_{0}^{1}x^{m}(1-x)^{n}\,dx=\int_{0}^{1}x^{n}(1-x)^{m}\,dx$ and hence evaluate $\displaystyle\int_{0}^{1}x(1-x)^{10}\,dx$.}
\q{218}{Evaluate the following definit integral: $\displaystyle\int_{0}^{1}\dfrac{x^{2}}{\sqrt{1-x}}\,dx$}
\q{219}{Let f(x) be a function satisfying $f(x)+f(-x)=x^{2}$ for all x. Evaluate the definite integral $\displaystyle\int_{-2}^{2}f(x)\,dx$}
\q{220}{Let $\displaystyle\int_{0}^{3}f(x)\,dx=20,\ \int_{0}^{10}f(x)\,dx=45$. Evaluate $\displaystyle\int_{3}^{10}f(x)\,dx$.}
\q{221}{Given that $\displaystyle\int_{0}^{9}f(x)\,dx=16$ and $\displaystyle\int_{0}^{3}f(x)\,dx=12$. Evaluate $\displaystyle\int_{3}^{9}(2f(x)+5)\,dx$.}
\q{222}{If $f(4-x)=f(x)$, show that $\displaystyle\int_{0}^{2}f(2x)\,dx=\int_{0}^{2}f(x)\,dx$}
\q{223}{Show that $\displaystyle\int_{-1}^{1}(1-x)f(1-x^{2})\,dx=\int_{-1}^{1}f(1-x^{2})\,dx$.}
\q{224}{If $f(a+x)+f(a-x)=0$, show that $\displaystyle\int_{0}^{2a}f(x)\,dx=0$}
\q{225}{Show that $\displaystyle\int_{-a}^{a}\dfrac{f(x)}{1+e^{x}}\,dx=\int_{0}^{a}f(x)\,dx$ where f(x) is an function. Hence or otherwise evaluate $\displaystyle\int_{-a}^{a}\dfrac{x^{2}}{1+e^{x}}\,dx$.}
\q{226}{Given the function f(x) satisfies the following condition: $f(x+2)=f(x)$ Evaluate: $\displaystyle\int_{-1}^{5}f(x)\,dx$.}
\q{227}{Show that: $\displaystyle\int_{0}^{\pi}\dfrac{\sin x\cos x}{1+e^{\sin x}}\,dx=0$}
\q{228}{Show that: $\displaystyle\int_{0}^{\pi}x(\pi-x)\cos(\pi-x)\,dx=0$}
\q{229}{Show that: $\displaystyle\int_{1}^{3}\dfrac{(x-2)^{3}+4(x-2)}{1+e^{(x-2)^{2}}}\,dx=0$}
\q{230}{Evaluate the following definite integral: $\displaystyle\int_{1}^{3}\dfrac{x^{2}}{x^{2}+(4-x)^{2}}\,dx$}
\q{231}{Evaluate the following definite integral: $\displaystyle\int_{0}^{\ln2}\dfrac{e^{2}}{e^{x}+2e^{-x}}\,dx$}
\q{232}{Evaluate the following definite integral: $\displaystyle\int_{1}^{3}\dfrac{\ln(x+2)}{\ln(x-6)+\ln(x+2)}\,dx$}
\q{233}{Evaluate the following definite integral: $\displaystyle\int_{0}^{\pi}\dfrac{e^{\sin x}-e^{\cos x}}{e^{\sin}+e^{\cos x}}(\sin x+\cos x)\,dx$}
\q{234}{Evaluate: $\displaystyle\int_{0}^{\sqrt{7}}\dfrac{x}{\sqrt{x^{2}+9}}\,dx$}
\q{235}{If $f(x)=x(\pi)=0$ , show that $\displaystyle\int_{0}^{\pi}f(x)\sin x\,dx=-\int_{0}^{\pi}f''(x)\sin x\,dx$}
\q{236}{Lef f(x) = $\begin{cases}x^{2}&\text{if }0\leq x<2\\4-x&\text{if }2\leq x\leq4\end{cases}$ Compute : $\displaystyle\int_{0}^{4}f(x)\,dx$}
\q{237}{Evaluate: $\displaystyle\int_{-6}^{4}f(x)\,dx$ where $f(x)=|x-1|+|x+2|$}
\q{238}{Evaluate: $\displaystyle\int_{0}^{1}x^{2}(1-x)^{\frac{3}{2}}\,dx$}
\q{239}{Show that $\displaystyle\int_{\pi/5}^{3\pi/10}\dfrac{\sqrt{\sin x}-\sqrt{\cos x}}{\sqrt{2}+\sin^{2}x\cos^{2}x}\,dx=0$}
\q{240}{Given Show that $\displaystyle\int_{-3}^{3}\dfrac{f(x)}{1+e^{x}}\,dx=4$ and $\displaystyle\int_{-3}^{3}f(x)\,dx=10$. Evaluate $\displaystyle\int_{-3}^{3}\dfrac{f(x)}{1+e^{-x}}\,dx$.}
\q{241}{Evaluate: $\displaystyle\int_{-a}^{a}\dfrac{x^{3}\sin^{2}x}{x^{4}+2x^{2}+1}\,dx$}
\q{242}{Given $\displaystyle\int_{0}^{2}f(x)\,dx=4$ and $f(2-x)=f(x)=f(-x)$. Evaluate $\displaystyle\int_{1}^{3}f(x)\,dx$.}
\q{243}{Find the total area and the net signed area between the graph of $y=x^{2}-4$ and the x-axis over the interval (i) [-3,1] (ii) [-3,3]}
\q{244}{Sketch the region bounded by the given curves and find the area of the region: $y=2x,\ y=x^{2}$}
\q{245}{Sketch the region bounded by the given curves and find the area of the region: $y=\sqrt{x},y=x^{2}$}
\q{246}{Sketch the region bounded by the given curves and find the area of the region: $y=x^{4}-x^{2}\ ,\ y=1,\ y=1-x^{2}$}
\q{247}{Sketch the region bounded by the given curves and find the area of the region: $y=x^{2}$ and $y=|x|$}
\q{248}{Find the area bounded by the parabola $x=9-y^{2}$, and the y-axis.}
\q{249}{Find the ratio of the areas of the regions bounded by the parabola $y=x^{2}-4x+3$ and the coordinate axes.}
\q{250}{Find the area of the region bouded by the line y=(x-1)(x+1)(x-2) and the x-axis.}
\q{251}{Find the area of the region bounded by the line y=x+1 and the parabola $y^{2}=2x+10$.}
\q{252}{The demand function is $p=30-2x-x^{2}$ and supply function is $p=-6+4x+x^{2}$. Find consumer surplus and producer surplus.}
\q{253}{At equilibrium, th quantity sold is $x_{e}=4$ and the price is $p_{e}=16$ . Consumer surplus is 16 and producer surplus is 8. Both the demand and supply curves are lineqar. Find the demand function D(x) and the supply function S(x).}
\q{254}{The supply function is p=ln(x+2) and the quantity sold is $x_{e}=2$. sonsumer surplus is 6-2 ln 2. Find Equilibrium price.}
\q{255}{The supply function is p=ln(x+2) and the quantity sold is $x_{e}=2$. sonsumer surplus is 6-2 ln 2. Linear demand function p=a-bx.}
\q{256}{If the demand function for a particular commodity is p=100-2x and the market price is $p_{e}=Rs.64$, find the consumer surplus using two methods. Integrate with respect to x.}
\q{257}{If the demand function for a particular commodity is p=100-2x and the market price is $p_{e}=Rs.64$, find the consumer surplus using two methods. Integrate with respect to p.}
\q{258}{Find the volumes of the solids obtained by rotating the region bounded by the given curves about the indicated axis: $y=x^{2},x=2,y=0;\ x-axis$}
\q{259}{Find the volumes of the solids obtained by rotating the region bounded by the given curves about the indicated axis: $y=\sqrt{x-5},x=6,x=8,y=0;\ x-axis$}
\q{260}{Find the volumes of the solids obtained by rotating the region bounded by the given curves about the indicated axis: $y=x^{2}-2x,y=0;\ x-axis$}
\q{261}{Find the area bounded by the curves $y=13-x^{2}$ and y=2x+5.}
\q{262}{Find the area of the region bounded by the parabolas $y=4-x^{2}$ and y=x$^{2}$}
\q{263}{The supply curve for a good is $p_{s}=\sqrt{16+x}$ . The demand curve is linear with vertical intercept 73 and passes through the equilibrium point. If the quantity sold is $x_{e}=33$, find the equation of the demand curve p$_{d}$}
\q{264}{If the supply function for a particular commodity is 25p=(x+10)$^{2}$ and the market price is Rs 64, find the producer surplus by two methods: Integrate with respect to x.}
\q{265}{If the supply function for a particular commodity is 25p=(x+10)$^{2}$ and the market price is Rs 64, find the producer surplus by two methods: Integrate with respect to p.}
\q{266}{A solid has a circular base of radius 2. Parallel cross-sections perpendicular to the base are equilateral triangles. Find the volume of the solid.}
\q{267}{Find the volume of the solid obtained by rotating the region bounded by $y=x^{3},y=64$ and $x=0$ about the y-axis.}
\q{268}{A cone is formed by rotating the straight line y=2x around the x-axis from x=0 to x=5 cm. Find the volume of the cone.}
\q{269}{A mug is designed by rotating the line y=3x+1 around the x-axis from x=0 to x=4 cm. Find the volume of the mug.}
\q{270}{Find the moment of inertia I of a uniform rectangular plate ABCD about the edge AD when AD=3, AB=8, and the mass per unit area is 4.}
\q{271}{A particle moves along a straight line with acceleration $a(t)=12t-36(m/s^{2})$ and initial velocity v(0)=30 m/s. Find the velocity function v(t).}
\q{272}{A particle moves along a straight line with acceleration $a(t)=12t-36(m/s^{2})$ and initial velocity v(0)=30 m/s. Find the displacement of he particle during the time period $1\leq t\leq6$}
\q{273}{A particle moves along a straight line with acceleration $a(t)=12t-36(m/s^{2})$ and initial velocity v(0)=30 m/s. Find the distance traveled during this time period.}
\q{274}{A particle moving along a coordinate line at time t=0 is at a position 2cm from the origin and traveling at a velocity of 5 cm/sec. If the acceleration of the particle is given by $a(t)=3-4(t+2)^{-3}$. Find the velocity and the position of the particle as function of t.}
\q{275}{A projectile is launched vertically upward from an initial height of $s_{o}$ feet with an initial velocity of $v_{o}ft/s$. Neglecting air resistance, show that its height s(t) above the ground at time t is given by $s(t)=-\dfrac{1}{2}gt^{2}+v_{o}t+s_{o}$ where g>0 is the magnitude of the acceleration due to gravity.}
\q{276}{In a bacterial culture, the growth rate is proportional to the number of bacteria present. Initially there are 200 bacteria. After 2 hours, there are 1800 bacteria. Find the number of bacteria at time t hours.}
\q{277}{In a bacterial culture, the growth rate is proportional to the number of bacteria present. Initially there are 200 bacteria. After 2 hours, there are 1800 bacteria. How many bacteria will there be after 4 hours?}
\q{278}{The rate of population growth of a small city (in people per year) is modeled by $P'(t)=800+60t$, where t is the number of years after 2025. if the population in 2025 (t=0) was 25,000 people: Find the population function P(t).}
\q{279}{The rate of population growth of a small city (in people per year) is modeled by $P'(t)=800+60t$, where t is the number of years after 2025. if the population in 2025 (t=0) was 25,000 people: Estimate the population in 2023.}
\q{280}{The rate of population growth of a small city (in people per year) is modeled by $P'(t)=800+60t$, where t is the number of years after 2025. if the population in 2025 (t=0) was 25,000 people: Calculate the total population growth from 2028 to 2033.}
\q{281}{Let f be continuous and suppose: $\displaystyle\int_{0}^{2}f(x)\,dx=5$. If the average value of f on [0,4] is twice the average value on [0,2], find $\displaystyle\int_{2}^{4}f(x)\,dx$.}
\q{282}{Let f be continuous and suppose $\displaystyle\int_{1}^{3}f(x)\,dx=8$. If the average value f on [1,4] is three times the average on [3,4], find $\displaystyle\int_{3}^{4}f(x)\,dx$.}
\q{283}{Let f be continuous and suppose $\displaystyle\int_{0}^{2}f(x)\,dx=10$ and $\displaystyle\int_{2}^{4}f(x)\,dx=6$. If the average value of f on [0,3] equals the average value on [3,4], find $\displaystyle\int_{1}^{4}f(x)\,dx$}
\q{284}{A patient receives an IV bolus dose of a drug. The concentration in the blood is modeled by $Ct=12e^{-0.3t}(mg/L)$ where t is the time in hours since that does was taken. the minimum effective concentration (MEC) is 3 mg/L. The drug is considered effective if total AUC in a 24-hour period is at least 150 mg-h/L. Find the time T (in hours) when the concentration reaches the MEC.}
\q{285}{A patient receives an IV bolus dose of a drug. The concentration in the blood is modeled by $Ct=12e^{-0.3t}(mg/L)$ where t is the time in hours since that does was taken. the minimum effective concentration (MEC) is 3 mg/L. The drug is considered effective if total AUC in a 24-hour period is at least 150 mg-h/L. Calculate the AUC from t=0 to t=T for a single dose.}
\q{286}{A patient receives an IV bolus dose of a drug. The concentration in the blood is modeled by $Ct=12e^{-0.3t}(mg/L)$ where t is the time in hours since that does was taken. the minimum effective concentration (MEC) is 3 mg/L. The drug is considered effective if total AUC in a 24-hour period is at least 150 mg-h/L. Detemine if one dose provides enough total AUC over 24 hours. If not, how many doses (given at intervals of T hours) are needed in 24 hours to meet the requirement?}
\q{287}{Find the moment of inertia of a uniform rod of length 2a and mass M about a line through one end and perpendicular to the rod.}
\q{288}{Find the moment of inertia of a uniform rectangular manina about a line through the centre and parallel to a side.}
\q{289}{A particle moves along a straight line with acceleration $a(t)=6t-18(m/s^{2})$ and initial velocity v(0)=15 m/s. Find the velocity function v(t).}
\q{290}{A particle moves along a straight line with acceleration $a(t)=6t-18(m/s^{2})$ and initial velocity v(0)=15 m/s. Find the displacement of the particle during $0\leq t\leq3$.}
\q{291}{A particle moves along a straight line with acceleration $a(t)=6t-18(m/s^{2})$ and initial velocity v(0)=15 m/s. Find the total distance traveled during $0\leq t\leq3$.}
\q{292}{Growth rate of population: A town's population growth rate (in people per year) is given by P(t)=5000+200t where t is the number of year after 2020. If the population in 2020 (t=0) was 100,000. Find the population function P(t).}
\q{293}{Growth rate of population: A town's population growth rate (in people per year) is given by P(t)=5000+200t where t is the number of year after 2020. If the population in 2020 (t=0) was 100,000. What will be the population be in 2030?}
\q{294}{Growth rate of population: A town's population growth rate (in people per year) is given by P(t)=5000+200t where t is the number of year after 2020. If the population in 2020 (t=0) was 100,000.Find the total population increase from 2022 to 2028.}
\q{295}{A temperature sensor moves slwoly along a 10-meter track, continuously measuring temperature. The temperature at position x meters is given by $T(x)=20+x$, where $T(x)$ is in $^{o}C$ Find the average temperature measured by the sensor as it moves from x=0 to x=10 meters.}
\q{296}{A patient is given an intravenous dose oof an antibiotic. The concentration in the bloodstream over time is modeled by: $C(t)=8e^{-0.15t}(mg/L)$ where t is an hours. Find the total drug exposure (AUC) from t=0 to t=6 hours.}
\q{297}{Let A denote the amount of money deposited in a bank at time t, which grows over time with constant compound interest rate r>0. Construct a differential equation which represent the growth of amount over time. (hint: the rate of change of amount with respect to time is proportional to amount itself).}
\q{298}{Let T denote the temperature of a body at time t. Suppose that the effect of the surrounding medium is negligible. Construct a differential equation that represent the temperature of the body at any time t, assuming the body is losing hat due to some internal processes.}
\q{299}{Let N denote the amount of radioactive substance present at time t. Construct a differential equation that models the radioactive decay of the substance. Hint: (radioactive decay is a natural process by which an unstable atoms nucleus loses energy by emitting radiation. The rate at which a radioactive substance decays is proportional to the amount of the substance present at any given time).}
\q{300}{For an RC circuit, recall the differential equation representing the charging of a capacitor. Deduce the differential equation that represents the discharging of the capacitor.}
\q{301}{For an RL circuit, recall the differenctial equation representing the charging of an inductor. Deduce the differential equation that represents the discharging of the inductor.}
\q{302}{Show that the following relation written against the differential equation is its solution. $x\tan y\dfrac{dy}{dx}-1=0\ ;\ x\cos y=b$}
\q{303}{Find the general solution of the differential equation: $\dfrac{dy}{dx}=\sqrt{y}\cos^{2}\!\sqrt{y}$}
\q{304}{Find the general solution of the differential equation: $\sqrt{x}\,\dfrac{dy}{dx}=e^{y+\sqrt{x}}$}
\q{305}{Find the general solution of the differential equation: $y-x\dfrac{dy}{dx}=2\left(y^{2}+\dfrac{dy}{dx}\right)$}
\q{306}{Find the general solution of the differential equation: $\sec^{2}x\tan y\,dx+\sec^{2}y\tan x\,dy=0$}
\q{307}{Find the general solution of the differential equation: $(e^{2x}+e^{-2x})dy+(e^{-2x}-e^{2x})dx=0$}
\q{308}{Find the general solution of the differential equation: $\dfrac{dy}{dx}=\cos^{2}\!\left(\dfrac{y}{x}\right)+\dfrac{y}{x}$}
\q{309}{Find the general solution of the differential equation: $\dfrac{dy}{dx}+\dfrac{2xy}{2y+1}=x$}
\q{310}{Solve the initial value problem: $\left(y+\sqrt{x^{2}+y^{2}}\right)dx-x\,dy=0\ ;\ y=0$ when $x=1$}
\q{311}{Show that the relation $\tan y=10\,(e^{x}-1)^{2}$ represents a particular solution for the differential equation $2e^{x}\tan y+(1-e^{x})\sec^{2}y\dfrac{dy}{dx}=0$}
\q{312}{Solve the differential equation $2e^{x}\tan y\,dx+(1-e^{x})\sec^{2}y\,dy=0\ ;\ y\in\left(-\dfrac{\pi}{2},\dfrac{\pi}{2}\right)y\neq0$}
\q{313}{Find the general solution of the differential equation $\dfrac{dy}{dx}=\dfrac{x+y}{y-x}$}
\q{314}{The slope of tangent at any point to a family of curves $f(x,y,c)=0$ is given by the rule $m(x,y)=\dfrac{x^{2}+3y^{2}}{2xy}$ x and y are not equal to zero. Find particular memvers of the family passing through the point $P_{o}(2,6)$}
\q{315}{In a culture, there were 1000 bacteria initially and were found to be 900 after five hours. Use differential equation to find the number of bacteria after 7 hours.}
\q{316}{The number of bacteria in a culture currently is 1000. Three hours ago, it was 750. What was the number of bacteria five hour ago?}
\q{317}{A radioactive isotopes has a half-life of 12 hours. The initial mass is 160 grams. How much mass will remain after 36 hours?}
\q{318}{A radioactive isotopes has a half-life of 12 hours. The initial mass is 160 grams. After how many hours will only 10 grams remains?}
\q{319}{Suppose the inflation rate for the cost of a basket of goods is 18\% per year. If the current price level is Rs. 1000, use differential equation to find the price level after 5 years.}
\q{320}{Use differential equation to find the current price a good if the deflation rate is 11\% per year and its price after tne years is experted to be Rs 5000.}
\q{321}{Aslam deposited Rs 5000 on opening an account in a bank 15 years ago. Using differential equation, find the constant annual rate of growth if he closes his acount now by reciving Rs 15,000.}
\q{322}{A can of soda at $4^{o}C$ is placed in a room at $24^{o}C$. After 30 minutes, its temperature is $12^{o}C$. Find the temperature after 1 hour.}
\q{323}{An RL circuit has a battery of 12 volts, a resistor of 60Ohms, and an inductor of 2 Henries. If the circuit is closed at time t=0 with initial current zero, find the current I(t) flowing in the circuit at any time t.}
\q{324}{An RC circuit has an emf of 12 volts, resistance $R=6\Omega$ and capacitance C=3 farad. Initially, the capacitor is uncharged. Find the current I(t) at any time t.}
\q{325}{In a culture, there are 100 bacteria initially and are doubled and are doubled in 2 hours. Find the number of bacteria 7 hours later.}
\q{326}{A certain radiactive sybstance has a mass of 80 grams initially. After 10 hours, only 50 grams remain.  Find decay constant.}
\q{327}{A certain radiactive sybstance has a mass of 80 grams initially. After 10 hours, only 50 grams remain. How much of the substance will remain after 24 hours?}
\q{328}{A certain radiactive sybstance has a mass of 80 grams initially. After 10 hours, only 50 grams remain. After how many hours will only 20 grams remain?}
\q{329}{Suppose an economy is experiencing deflation at 4\% per year. If the current price of an item is Rs 500, find its expected price after 3 years.}
\q{330}{A cup of coffee is served at $90^{o}C$ in a room where the temperature is $25^{o}C$. After 10 minutes the coffee cools to $60^{o}C$. Find the temperature of the coffee after 20 minutes.}
\q{331}{An RL circuit has an emf of 5 volts, a resistance of 50 ohms, an inductor of 1 henry and no initial current. Find the current in the circuit at any time t.}
\q{332}{A(-3,-7), B(4,5) and C(-1,2) are the vertices of $\triangle ABC$. Find the equation of right bisector of side BC.}
\q{333}{The vertices of $\triangle PQR$ are P(1,2), Q(5,-6) and R(7,-2). Find the equation of all three medians.}
\q{334}{The vertices of a triangle are A(2,-1), B(4,3) and C(6,-1). Find the equations of altitudes and show that they are concurrent.}
\q{335}{Find the equations of medians, altitudes and right bisectors of the trangle whose vertices are A(2,5), B(-3,4) and C(6,-8).}
\q{336}{Convert the following equation into: 2x+3y-6=0 (i) Slope intercept form (ii) Two point form (iii) Intercept form (iv) Point slope form (v) Normal form (vi) Symmetric form.}
\q{337}{Convert the following equation into: x-2y+4=0 (i) Slope intercept form (ii) Two point form (iii) Intercept form (iv) Point slope form (v) Normal form (vi) Symmetric form.}
\q{338}{Convert the following equation into: 3x+4y-10=0 (i) Slope intercept form (ii) Two point form (iii) Intercept form (iv) Point slope form (v) Normal form (vi) Symmetric form.}
\q{339}{Find the equation of family of lines through the point of intersection of the lines x+2y-3=0, 3x-y+4=0. find the member of the family which is: (i) parallel to a line with slope $\dfrac{3}{2}$ (ii) Perpendicular to the line 3x+4y=0}
\q{340}{Find the family of lines through the point of intersection of the lines 2x-3y+5=0, x+4y-7=0. Find the member of the family which is (i) parallel to y=2x+1 (ii) perpendicular to y=x.}
\q{341}{Check whether the following lines are concurrent or not. If concurrent, then find the point of concurrency x+2y-7=0; 3x-y-4=0; 2x+35y-145=0.}
\q{342}{Find the condition that the lines $y=m_{1}x+c_{1}\ ;\ y=m_{2}x+c_{2}$ and $y=m_{3}x+c_{3}$ are concurrent.}
\q{343}{Transform the equation 3x-10y+13=0 into: (i) slope intercept form (ii) Two intercept form (iii) Normal form (iv) point slope form (v) Two point form (vi) Symmetric form.}
\q{344}{Triangle ABC has vertices A(2,3), B(4,5) and C(-2,7). Find the equations of medians of the triangle ABC.}
\q{345}{A(1,2), B(5,-1) and C(-2,-4) are the vertices of $\triangle ABC$. Find the equations of altitudes of triangle ABC.}
\q{346}{Triangle ABC has vertices A(-3,2), B(5,4) and C(3,-8). Find the equations of right bisectors of triangle ABC.}
\q{347}{Check whether the following lines are concurrent or not. If concurrent then find the point of concurrency: 3x-2y-2=0 , x+y-4=0 , 2x-y-2=0}
\q{348}{Find the family of lines through the point of intersection of the lines 2x-3y-14=0 and 2x+y-10=0.}
\q{349}{Find the family of lines through the point of intersection of the lines 2x - 3y - 14 = 0 and 2x + y - 10 = 0\par Find the member of the family which is:\par (a) Parallel to the line with slope $-\dfrac{1}{2}$\par (b) perpendicular to the line 2x - 3y + 2 = 0}
\q{350}{Find the angle measured from the line $\ell_{1}$ to the line $\ell_{2}$, where. $\ell_{2}=$ joining $(-7,2)$ and $(1,3)$. Also find the acute angle.}
\q{351}{Find the angle measured from the line $\ell_{1}$ to the line $\ell_{2}$, where. $\ell_{2}=$ joining $(2,-1)$ and $(4,3)$. Also find the acute angle.}
\q{352}{The angle from the line ax-y+2=0 to the line x+y-4=0 is $30^{o}$. Find the value of a.}
\q{353}{Find the interior angles of the triangle whose vertices are: A(2,3), B(6,3), C(4,7).}
\q{354}{Find the interior angles of the triangle whose vertices are:L(1,2), M(5,5), N(7,2)}
\q{355}{Find the interior angles of the quadrilateral whose vertices are A(2,1), B(3,1), C(4,5) and D(-1,5).}
\q{356}{Find the interior angles of the kite whose vertices are A(0,0), B(2,3), C(4,0) an D(2,-3).}
\q{357}{Find the area of the triangle formed by the points A(-3,2), B(4,-1), and C(1,5).}
\q{358}{By computing the area of triangle ABC check wheter the points are collinear where A(0,4), B(3,2) and C(6,0).}
\q{359}{The coordinates of the points are A(0,$\lambda$), B(-2,1) and C(-3,-2). By computing the area bounded by ABC, dins the value of $\lambda$ if points are collinear.}
\q{360}{Find the angle of the triangle whose vertices are A(1,2), B(7,3) and C(4,9).}
\q{361}{Find the value of k such that the lines represented by the equation: $(k-8)x^{2}-9xy+3ky^{2}$ are perpendicular to each other.}
\q{362}{Prove that the lines represented by the equation $9y^{2}+6xy+x^{2}=0$ are coincident.}
\q{363}{Find a joint equation of the lines through the origin and perpendicular to the lines $3x^{2}+5xy-2y^{2}=0$.}
\q{364}{The paths of two rays froma star are represented by the equation $2x^{2}+5xy-3y^{2}=0$. Find the equation of the two lines and determine the angle between the rays.}
\q{365}{The trajectories of two projectiles in a physics experiment are modeled by the homogeneous equation $9x^{2}-12xy+4y^{2}=0$. Prove that this represents two coincident lines and angle between them.}
\q{366}{The stresses in a metal plate are represented by the homogeneous equation $3x^{2}-8xy-3y^{2}=0$. Show that this equation represents a pair of straight lines through the origin.\par Find: (i) the equations of lines represented by this equation. (ii) angle between them.}
\q{367}{Find the equation of lines represented by the following and also find measure of the angle between them: $7x^{2}-8xy+y^{2}=0$}
\q{368}{Find the equation of circle which passes through the origin and cuts off intercepts 2 and 3 on the x-axis and y-axis respectively.}
\q{369}{Find the equation of a circle passinf through the points (1,3), (2,-1) and having its centre on the line x+y=3. Does the point (4,1) lie on the circle?}
\q{370}{Find the equation of the circle passing through the point (-3,-4) and touching the line 3x-y-1=0 at the point (1,2).}
\q{371}{Find the equation of the circle passing through the point of intersection of the lines x-y-5=0, 2x+y-4=0 and it has its centre at (5,0).}
\q{372}{Find the equation of the circle touching the u-axis and also the straight line 3x-4y=0 at the point (8,6).}
\q{373}{Find the equation of the circle passing through the points A(3,2) and B(3,-2) and touching the line x+2y+3=0}
\q{374}{Prove that the locus of the point of intersection of the lines $x\cos\theta+y\sin\theta=\alpha,\ s\sin\theta-y\cos\theta=b$. is a circle with centre at (0,0) and radius $\sqrt{a^{2}+b^{2}}$ , where (equation) is a parameter.}
\q{375}{Show that the equation $5x^{2}+5y^{2}+20x+25y+10=0$ represents a circle. Also find its centre and radius.}
\q{376}{A point moves so that the sum of the squares of its distance from two points A(1,1) and B(2,2) is equal to 4. Prove that its locus is a circle.}
\q{377}{Find the equation of the circle passing through the points (-3,3), (-1,5) and (5,1).}
\q{378}{A circle with its centre on the y-axis passes through the origin and the point (a,b), where b>0. Find its equation.}
\q{379}{Find the equation of the circle passing through (1,2) and touching the line x-2y+4=0 at (2,3).}
\q{380}{Find the equations of the circles passing through the points A(3,4) and B(3,-4) and touching the line x+2y+9=0.}
\q{381}{Find the equation of the circle which touches the y-axis at a distance of 3 units from the origin and cuts off an intercept of 8 units on the x-axis.}
\q{382}{Find the equation of the tangents to the circle $x^{2}+y^{2}=4$: which are inclined at an angle of $45^{o}$ to the axis of x.}
\q{383}{Find the equation of the tangents to the circle $x^{2}+y^{2}=4$: which are parallel to the line 3x+4y+1=0.}
\q{384}{Find the equations of the tangents to the circle $x^{2}+y^{2}-4x-2y-4=0$ which are parallel to the line 3x-4y+2=0.}
\q{385}{Find the equations of the tangents drawn from: the point (-3,7) to the circle $x^{2}+y^{2}=8$. Also find the points of contact.}
\q{386}{Find the equations of the tangents drawn from: the point (12,-4) to the circle $x^{2}+y^{2}-4x-2y-20=0$. Also points of contact.}
\q{387}{A circle has the parametric equation: $x=3+5\cos\theta.\ y=-1+5\sin\theta$ Find the Cartesian equation of the circle.}
\q{388}{A circle has the parametric equation: $x=3+5\cos\theta.\ y=-1+5\sin\theta$ Find the centre and radius of the circle.}
\q{389}{A circle has the parametric equation: $x=3+5\cos\theta.\ y=-1+5\sin\theta$ Find the parametric coordinates of the point P on the circle corresponding to $\theta=\dfrac{\pi}{2}$.}
\q{390}{A circle has the parametric equation: $x=3+5\cos\theta.\ y=-1+5\sin\theta$ Find the equation of the tangent to the circle at point P.}
\q{391}{A circle has the parametric equation: $x=3+5\cos\theta.\ y=-1+5\sin\theta$ Find the equation of the normal to the circle at point P.}
\q{392}{A circle has the parametric equation: $x=3+5\cos\theta.\ y=-1+5\sin\theta$ Find the equation of the tangent to the circle at the point where $\theta=\pi$.}
\q{393}{Find the equation of the tangent and normal to the circle $x^{2}+y^{2}-4x+6y+12=0$ at the point on the circle whose ordinate is -2.}
\q{394}{Tangents are drawn from (4,5) to the circle $x^{2}+y^{2}=29$. Find an equation of the line joining the points of contanct (The line is called the chord of contact).}
\q{395}{Find a combined equation of the pair of tangents drawn from (5,0) to the circle $x^{2}+y^{2}=9$.}
\q{396}{Find the equations of the pair of tangents drawn from (4,-1) to the circle $x^{2}+y^{2}-2x+4y=0$. Also find the points of contact.}
\q{397}{Find the length of intercept cut off from the line y=mx+c by the circle $x^{2}+y^{2}=a^{2}$.}
\q{398}{Find the focus, vertex and directrix of the following parabola: $x^{2}-6x-12y+9=0$.}
\q{399}{Write the equation of the parabola with given elements: Focus (-3,2); directrix x=5.}
\q{400}{Show that the equation of a parabola with focus at $(\alpha\cos\theta,\ \alpha\sin\theta)$ and directrix $x\cos\theta+y\sin\theta+\alpha=0$ is $(x\sin\theta-y\cos\theta)^{2}=4\alpha(x\cos\theta+y\sin\theta)$.}
\q{401}{Prove that the chord of contact of the tangents drawn from any point on the directrix of the parabola $y^{2}=4ax$ always passes through the focus.}
\q{402}{The parabola $y=ax^{2}+bx+c$ has a vertex at (3,-2) and passes through the point (1,6). Find the values of a,b and c.}
\q{403}{Find the equation of the tangent and normal to the parabola: $x^{2}-6x-2y+11=0$ at the point (5,3).}
\q{404}{Find the equation of the tangent to the parabola $y^{2}=12x$ which is parallel to the line 3x-y+4=0. Also find the point to tangency.}
\q{405}{Analyze the parabola $x^{2}-4x-4y-2=0$ and sketch its graph.}
\q{406}{Find the equations of the tangent and normal to the parabola $y^{2}=4ax$ at the point $(x_{1},y_{1})$.}
\q{407}{Find an equation of the ellipse with given data: Centre (-3,1), major axis parallel to the x-axis and of length 10 units, minor axis parallel to the y-axis and of the length 6 units.}
\q{408}{Find the centre, foci, eccentricity, vertices and directrices of the ellipse whose equation is given : $4x^{2}+9y^{2}=36$}
\q{409}{Find the centre, foci, eccentricity, vertices and directrices of the ellipse whose equation is given : $4x^{2}+24x+9y^{2}-36y+36=0$}
\q{410}{Find the equation of the ellipse as the locus of points P(x,y) such that the sum of the distance from P to the points (-3,0) and (3,0) is 10.}
\q{411}{What are the dimensions of the smallest rectangle aligned to the axes that completely contains the ellipse $\dfrac{x^{2}}{16}+\dfrac{y^{2}}{9}=1$?}
\q{412}{Find the dimensions of the rectangle with greates perimeter that can be inscribed in the ellipse $\dfrac{x^{2}}{16}+\dfrac{y^{2}}{9}=1$}
\q{413}{Find the value of k for which the line y=2k+k is tangent to the ellipse $\dfrac{x^{2}}{9}+y^{2}=1$.}
\q{414}{Determine whether the line y=2x-3 is a tangent, secant or external to the ellipse $\dfrac{x^{2}}{9}+\dfrac{y^{2}}{4}=1$.}
\q{415}{Find the equation of an ellipse with centre at (0,0), a horizontal major axis of length 8, and that touches the line y=x+5.}
\q{416}{Show that the family of lines given by $x\cos\theta+y\sin\theta=\sqrt{a^{2}\cos^{2}\theta+b^{2}\sin^{2}\theta}$ where $\theta$ is a parameter and a>b are all tangent to the ellipse $\dfrac{x^{2}}{a^{2}}+\dfrac{y^{2}}{b^{2}}=1$}
\q{417}{If tangent to the ellipse $\dfrac{x^{2}}{a^{2}}+\dfrac{y^{2}}{b^{2}}=1$ with centre C meets the major and minor axes at P and Q respectively. Provethat $\dfrac{a^{2}}{|CP|^{2}}+\dfrac{b^{2}}{|CQ|^{2}}=1$.}
\q{418}{Find an equation of the ellipse with foci (2,1), (2,7) major axis of length 10 and sketch its graph.}
\q{419}{Find an equation of the ellipse with centre (0,0), focus (5,0) and vertex (6,0). Sketch its graph.}
\q{420}{Find the centre, foci, eccentricity, vertices and directrices of the ellipse whose equation is $16x^{2}+9y^{2}-64y+54y+1=0$ . Also sketch its graph.}
\q{421}{Find the equation of the tnagent and normal to the ellipse $\dfrac{x^{2}}{a^{2}}+\dfrac{y^{2}}{b^{2}}=1$ at the point (x$_{1}$ y$_{1}$)}
\q{422}{A hyperbola has its centre at the origin, transverse axis of length 14 along the y-axis and conjugate axis of length 10 along the x-axis. Find the equation of the hyperbola.}
\q{423}{A hyperbola has a horizontal transverse axis of length 8. One of its foci is at the point $(1+2\sqrt{5},3)$ and its eccentricity is $\dfrac{\sqrt{5}}{2}$ . Find the equation of the hyperbola.}
\q{424}{A hyperbola has a vertical transverse axis of length 8, centre at (2,-1) and eccentricity $\dfrac{\sqrt{13}}{2}$ . Find the equation of the hyperbola.}
\q{425}{A hyperbola has its centre at (-1,4) a focus at (5,4) and the corresponding directrix given by the line x=-3. Find the equation of the hyperbola.}
\q{426}{A hyperbola has a vertical transverse axis, exxentricity $\dfrac{\sqrt{19}}{4}$ , and latus rectum of length $\dfrac{3}{2}$ . The hyperbola passes through the points (2,8) and (2,0). Find the equation of the hyperbola.}
\q{427}{The foci of a hyperbola coincide with the foci of the ellipse $\dfrac{x^{2}}{25}+\dfrac{y^{2}}{9}=1$ . If the eccentricity of the hyperbola is 2. Find its equation.}
\q{428}{Find the equations of the tangents to the hyperbola $\dfrac{x^{2}}{4}-\dfrac{y^{2}}{12}=1$ which are parallel to the line 2x-y=5.}
\q{429}{A hyperbola for which a=b is called rectangular. Show that a hyperbola is rectangular if and only if its asymptotes are perpendicular to one another.}
\q{430}{Find the equations of the tangents to the hyperbola $7x^{2}-3y^{2}=105$ that pass through the point (3,1).}
\q{431}{Find the condition that the line Ix+my=1 is a tangent to the hyperbola $\dfrac{x^{2}}{a^{2}}-\dfrac{y^{2}}{b^{2}}=1$.}
\q{432}{A hyperbola has a horizontal transverse axis of length 12, centre at (-2,5) and eccentricity $\dfrac{5}{3}$ . Find the equation of the hyperbola.}
\q{433}{A hyperbola has a horizontal trnsverse axis, eccentricity $\dfrac{\sqrt{13}}{3}$ and latus rectum of length $\dfrac{8}{3}$ . The hyperbola passes through the points (8,2) and (2,2) Find the equation of the hyperbola.}
\q{434}{Find an equation of the hyperbola with given data and sketch the graph. (i) Foci: $(-3,5\pm5\sqrt{2})$ (ii) Length of transverse axis=6}
\q{435}{Find an equation of the hyperbola with given data and sketch the graph (i) Foci (4,-3), (4,7) and (ii) one vertex (4,5).}
\q{436}{Find the centre, foci, eccentricity, vertices and equation of the directrices of the hyperbola. $9x^{2}-36x-4y^{2}+8y-4=0$. Also sketch the graph.}
\q{437}{Find equation of the tangent and normal to the hyperbola $\dfrac{x^{2}}{a^{2}}-\dfrac{y^{2}}{b^{2}}=1$ at the point $(x_{1},y_{1})$.}
\q{438}{Find the equations of the tangent to the hyperbola $2x^{2}-3y^{2}=24$ that pass through the point (2,-4).}
\q{439}{A coment has a parabolic orbit with the Sum at the focus. When the comet is 100 millions km from the sum. the line joining the Sun and the comet makes an angle of $60^{o}$ with the axis of the parabola. How close will the comet get to the sun?}
\q{440}{Find an equation of the parabola formed by the cables of suspension bridge whose span is $\ell$m and the vertical height of the supporting towers is hm.}
\q{441}{An athletic field features an elliptical running track. The track is 120 meters long along its major axis and 80 meters wide along its minor axis. (a) Find the distance from the centre to each focus of the ellipse. (b) If two water fountains are placed at the foci, how far apart are they?}
\q{442}{A landscaper wants to create an elliptical flower bed with a major axis of 10 meters and a minor axis of 6 meters. Using the gardener's method, two stakes are placed at the foci, and a loop of string is stretched around them to trace the ellipse: (a) How far apart should the two stakes be placed? (b) What should be the total length of the string loop?}
\q{443}{An architect designs a bridge with an elliptical arch. The arch has a span (width) of 24 meters and a height of 8 meters. The ellipse is centred at the origin with the major axis horizontal: (a) Write the equation of the ellipse representing the arch. (b) A boat needs to pass under the bridge. The boat is 6 meters wide. What is the maximum height of the boat that can pass through the arch if it must stay centred?}
\q{444}{In the Bohr-Sommerfeld moel, an electron moves in an elliptical orbit around a nucleus. For a hydrogen atom, the square of the orbital period T is proportional to the cube of the semi-major axis a: Consider an electron in an elliptical orbit with a semi-major axis $a=2.12\times10^{-10}$ and an eccentricity e=0.30. Th constant of proportionality for hydrogen is $k=2.96\times10^{-16}\,s^{2}/m^{3}$ (a) Determine the distance between the nucleus (at one focus) and the centre of the ellipse (b) Calculate the orbital period T of the electron. (c) Find the maximum and minimum distance of the electron from the nucleus during its orbit.}
\q{445}{LORAN (long sitance radio navigation) for aircraft and ships uses synchronized pulses transmitted by widely seprarted transmitting stations. These pulses travel at the speed of light (186,000 miles per second). The difference in the times of arrival of these pulses at an aircraft or ship is constant on a hyperbola having the transmitting stations as foci. Assume that two stations 50 miles apart, are positioned on a rectangular coordinates system at (-250,0) and (250,0) and that a ship is traveling on a path with coordinates (x,120) (a horizontal line 120 miles above the x-axis). Find the x-coorinates of the position of the ship when  the time difference between the pulses from the transmitting stations is 1200 microseconds (0.0012 second).}
\q{446}{A satellite orbits Earth in an elliptical path. The apogee (farthest distance from Earth) is A=42000 km, and the perigee (closest distance to Earth) is P=30000 km. (a) Find the semi-major axis a of the orbit.  (b) Find the distance from the centre to each focus c..  (c) Calculate the eccentricity e of the orbit.}
\q{447}{The end towers o a suspension bridge are 100 meters high and 540 meters apart. The suspension cable hangs the top of the towers in the form of a parabola and the lowest point of the cable is 25 meters above the horizontal road. Find the length of the supporting vertical cable at a horizontal distance of 81 meters from the centre of the bridge.}
\q{448}{An architectural arch is built in the shape of the upper half of an ellipse with semi-major axis 8 meters and semi-minor axies 5 meters. A horizontal beam is to be placed across the arch at a height where the width of the arch is exactly 10 meters. Calculate the exact height above ground level where this beam should be installed.}
\q{449}{A comet has a parabolic orbit with the Earth at the focus. When the comet s h km rom the earth, the line joining the comt and the Earth maker an angle of $30^{o}$ with the axis of the parabola. How close will the comet come to the Earth?}
\q{450}{In the Bohr-Sommerfeld model of the atoms, electrons move in elliptical orbits around the nucleus which is located at one focus. Consider an electron in an elliptical orbit with a semi-major axis $a=5.29\times10^{-11}m$ (the Bohr radius) and an eccentricity e=0.20. (a) Find the distance from the nucleus to the centre of the ellipse. (b) Determine the minimum distance (periapsis) and maximum distance (apoapsis) of the electron from the nucleus.}
\q{451}{The cable of a suspension bridge is parabolic. The roadway is 10 m. below the lowest point of the cable. The span of the bridge is 200 m, and the tops of the piers are 20 m above the roadway. Find the equation of the parabola. If the load is supported by vertical chains at intervals of 10m, find the length of the chain that is 20 m away from either pier.}
\q{452}{The moon revolves around Earth in an elliptical orbit, with the centre of Earth located at one focus. The major and minor axis of the orbit have lengths of 768,800 km and 767,600 km respectively. Find the greatest and least distance (the apogee and perigee) from the Earth's centre to the moon's centre.}
\q{453}{The apgee (the point in orbit farthest from Earth) and the perigee (the point in orbit closest to Earth) of an elliptical orbit of an Earth satellite are given by A and P respectively. Show that the eccentricity of the orbit is $e=\dfrac{A-P}{A+P}$.}
\q{454}{A cross-section of a parabolic reflector is shown in the figure\par\smallskip\includegraphics[height=4.2cm]{fg00.png}\par\smallskip The bulb is located at the focus and the opening at the focus is 20 cm. (a) Find an equation of the parabola.  (b) Find the diameter of the opening $|CD|$, 20 $cm$ from the vertex.}
\q{455}{Kamran covers 100 km in the first 2 hours, stops for 1 hour, and then drives 60 km in the next 3 hours. The displacement-time graph is shown in Figure\par\smallskip\includegraphics[height=4.2cm]{fg01.png}\par\smallskip (a) What is Kamran's average velocity?\par (b) What is his velocity during the last phase of the journey?}
\q{456}{The displacement-time graph of a moving car is shown in figure\par\smallskip\includegraphics[height=4.2cm]{fg02.png}\par\smallskip The car changes its direction during the journey, and its motion occurs in three distinct phases.\par (a) What is the car's velocity in each phase?\par (b) When did the car pass its starting point again?\par (c) What is the total distance travelled?\par (d) What is the average velocity?}
\q{457}{Figure\par\smallskip\includegraphics[height=4.2cm]{fg03.png}\par\smallskip is the velocity-time graph of Aslam's motion as he runs along a straight track. For the first 5 seconds, he accelerates uniformly from rest to a velocity of 8 m/s. He then maintains this velocity for a further 8 seconds. Find:  (a) The rate at which Aslam accelerates.  (b) The displacement from the starting point of Aslam after 13 seconds.}
\q{458}{A car is travelling along a straight road. At t=0 seconds, the car passes a point A with a velocity of 12 m/s, and this velocity is maintained up to t=35 seconds. The driver then applies the brakes, and the car decelerates uniformly, coming to rest at the point B at t = 40 seconds. (a) Sketch a velocity-time graph to illustrate the motion of the car.  (b) Find the displacement from A to B.}
\q{459}{Figure\par\smallskip\includegraphics[height=4cm]{fg04.png}\par\smallskip shows the velocity-time graph of an object moving in a straight line over a 12-second time interval.\par (a) At what time did the object change its direction of motion?\par (b) Calculate the total distance travelled and the displacement of the object during the 12 seconds.}
\q{460}{A train moving at 54 km/h is accelerated uniformly to 90 km/h in 20 s. Find the acceleration and the distance covered in this time.}
\q{461}{A particle moving in a straight line starts from rest and is accelerated uniformly to attain a velocity of 30 km/h in 5 seconds. Find the distance travelled by the particle in the last three seconds.}
\q{462}{A particle is projected vertically upwards with a velocity 100m / s and another is let fall from a height of 200 m at the same time. Find the height of the point where they meet each other.}
\q{463}{A ball is thrown vertically upward with a velocity of 20 m/s. Find, the maximum height reached, and the time taken to return to the thrower's hand.}
\q{464}{A rock is dropped from a cliff. One second later, another rock is dropped from the same height. (a) What is the distance between the two rocks 2 seconds after the first rock is dropped? (b) At what time will the distance between them be 14.7 m?}
\q{465}{The motion of a robotic arm moving along a straight track is described by the velocity function: $v(t)=6+t-t^{2}$ where v(t) is in meters per second and t is the time in seconds, valid for $0\leq t\leq5$. Find the distance travelled by the arm.}
\q{466}{Let (equation) be the position of a particle movine along s straight line at any time t, where x is in meters and t is in seconds. Find the velocity and the speed of the particle at any time t and at t= 1 seconds.}
\q{467}{A cyclist rides in a straight line for 20 minutes. He then stops and waits for 30 minutes, after which he return in a straight line tio the starting point in 10 minutes.\par\smallskip\includegraphics[height=3.6cm]{fg05.png}\par\smallskip shows the displacement time graph for his journey.\par (a) calculate the average velocity for each stage of the journey in km/h.\par (b) Calculate the average velocity for the whole journey. (c) Find the average speed for the whole journey.}
\q{468}{Javed drives from his home to Murree. He drives for 2 hours at an average velocity of 50 km/h. He then stop for lunch before continuting journey to Murree.\par\smallskip\includegraphics[height=4.2cm]{fg06.png}\par\smallskip shows a displacement time graph for his journey: (a) Work out the displacement of Murree from Javed's home.  (b) Work out his average velocity for his whole journey.}
\q{469}{\includegraphics[height=3.8cm]{fg07.png}\par\smallskip is the velocity time graph for a cyclist travelling in a straight line. (The time is in minutes and velocity is in meters per minute).\par (a) Calculate te total displacement of the cyclist from his starting position at the end of the 12 minute.\par (b) Determine the uniform deceleration during the final 4 minutes of motion.}
\q{470}{Ali drives a car 400 km in a straight line.\par\smallskip\includegraphics[height=4.6cm]{fg08.png}\par\smallskip is a sketch of the velocity time graph representing his motion. Ali starts with a velocity of u km/h and acceleratess to a velocity of 70 km/h in 3 hours. He then maintain his velocity for 3 hours, after which he slows down uniformly to rest over the next 1/2 hour. Find:\par (a) The value of u  (b) The acceleration of Ali during the first 3 hours of his drive.}
\q{471}{A stone is dropped from the top of a tower 45 m high. Find the time taken to reach the ground and the final velocity on impact.}
\q{472}{A ball is thrown vertically upward with a speed of 20 m/s. Find the maximum height attained an the time to reach the maximum height.}
\q{473}{Find the real roots of the following functions correct to two decimal places using the graphical method, within the specified interval of sign change. $f(x)=3x+\sin(x)-e^{x}[0,1]$}
\q{474}{Find the real roots of the following functions correct to two decimal places using the graphical method, within the specified interval of sign change. $f(x)=x^{4}-x-10\ [1,2]$}
\q{475}{Find the real roots of the following functions correct to two decimal places using the graphical method, by finding some interval of sign change. $f(x)=x^{3}-x-1$}
\q{476}{Find the real roots of the following functions correct to two decimal places using the graphical method, by finding some interval of sign change. $f(x)=xe^{x}-2$}
\q{477}{Find the real roots of the following functions correct to two decimal places using the graphical method, by finding some interval of sign change. $f(x)=\ln(x)-2$}
\q{478}{Show that the function $f(x)=x^{4}-200x-50$ has at least two real roots by finding its intervals of sign change.}
\q{479}{Find a root of the equation x-cos x=0 graphically, approximate to two decimal places.}
\q{480}{Find the real roots of the following functions correct to 3 decimal places using the bisection method, within the specified interval of sign change. $f(x)=x^{2}\sin(x)-1:[1,2]$}
\q{481}{Find the real roots of the following functions correct to 3 decimal places using the bisection method, within the specified interval of sign change. $f(x)=3x-\cos x-1:[0,1]$}
\q{482}{Find the real roots of the following functions correct to 3 decimal places using the bisection method, within the specified interval of sign change. $f(x)=x^{4}-2x-10:[1,2]$}
\q{483}{Find the real roots of the following functions correct to 3 decimal places using the bisection method, within the specified interval of sign change. $f(x)=x^{3}-3x-10:[2,3]$}
\q{484}{Find the real roots of nonlinear equations correct to 2 decimal places using the Regula -Falsi method, by finding some interval of sign change if not specified. $3x^{3}-9x^{2}+8=0\ ;\ [2,3]$}
\q{485}{Find the real roots of nonlinear equations correct to 2 decimal places using the Regula -Falsi method, by finding some interval of sign change if not specified. $x^{2}+4\sin x=0\ ;\ [-2.2,-1]$}
\q{486}{Find the real roots of nonlinear equations correct to 2 decimal places using the Regula -Falsi method, by finding some interval of sign change if not specified. $xe^{x}=2$}
\q{487}{Find the real roots of nonlinear equations correct to 2 decimal places using the Regula -Falsi method, by finding some interval of sign change if not specified. $3x+\sin x=e^{x}$}
\q{488}{Find the real roots of nonlinear equations using the Newton-Raphson method. $x^{3}+x^{2}=3(x+1)$ (x+1) near x=1 accurate to 2 decimal places.}
\q{489}{Find the real roots of nonlinear equations using the Newton-Raphson method. $e^{x}=4x$ near x = 2 accurate to 3 decimal places.}
\q{490}{Find the real roots of nonlinear equations using the Newton-Raphson method. $x^{3}+2x-5=0$ accurate to 4 decimal places.}
\q{491}{Find the real roots of nonlinear equations using the Newton-Raphson method. $3x+\sin x=e^{x}$ accurate to 3 decimal places.}
\q{492}{Find the real roots of nonlinear equations using the Newton-Raphson method. $\tan(X)=x$ around $x_{0}=4.3$ accurate to 4 decimal places.}
\q{493}{Use bisection method to find a root of $f(x)=x^{3}-3x-10$ within the interval [2,3] accurate up to two decimal places.}
\q{494}{Use Regula-Falsi method to find a root of $f(x)=x^{3}-3x-10$ within the interval [2, 3] accurate up to two decimal places.}
\q{495}{Use Newton-Raphson method to find a root of $f(x)=x^{3}+2x^{2}+10x-20$ correct two decimal places.}
\q{496}{Use Trapezoidal Rule to approximate the value of the following definite integrals. $\displaystyle\int_{0}^{1}e^{x}dx$ taking n=7}
\q{497}{Use Trapezoidal Rule to approximate the value of the following definite integrals. $\displaystyle\int_{1}^{7}\ln(x)\,dx$ taking n=6}
\q{498}{Use Trapezoidal Rule to approximate the value of the following definite integrals. $\displaystyle\int_{0}^{2}x^{3}dx$ taking n=8}
\q{499}{Use Trapezoidal Rule to approximate the value of the following definite integrals. $\displaystyle\int_{0}^{4}\sqrt{x}\,dx$ taking n=4}
\q{500}{Use Trapezoidal Rule to approximate the value of the following definite integrals. $\displaystyle\int_{-1}^{1}e^{-x^{2}}dx$ taking h = 0.4}
\q{501}{Use Trapezoidal Rule to approximate the value of the following definite integrals. $\displaystyle\int_{-1}^{1}(x^{4}-x+1)\,dx$ taking h = 0.2}
\q{502}{Use Trapezoidal Rule to approximate the value of the following definite integrals. $\displaystyle\int_{0}^{\pi}\sin(x^{2})\,dx$ taking $h=\dfrac{\pi}{6}$}
\q{503}{Show by the trapezium rule that the approximate area bounded by the curve $y=\sqrt{8-2x-x^{2}}$ and the x-axis using n=8 is 13.5 square units. Calculate the exact area by using the formula for finding the area of semi-circle.}
\q{504}{Setting h=0.5, use the trapezium rule to find the approximate area bounded by the curve $y=e^{x^{2}}$ from the line x=0 to the line x=4.}
\q{505}{Use trapezium rule to approximate the definite integral $\displaystyle\int_{0}^{1}e^{-x^{2}}dx$ by dividing the closed interval [0, 1] into 7 subintervals of equal length.}
\q{506}{Use Trapezium Rule to approximate the definite integral $\displaystyle\int_{0}^{6}\dfrac{x}{1+x^{2}}\,dx$ by subdividing the closed interval [0, 6] each of length 1. Also compare the answer by solving the integral using the Fundamental Theorem of Calculus.}
\q{507}{In a certain chemical reaction, the concentration of a substance at equilibrium is governed by the relation $x^{2}+0.2x-0.04=0$ , where x is the equilibrium concentration (in moles per liter). Use bisection method to find the positive equilibrium concentration of the substance accurate to two decimal places.}
\q{508}{The flow rate of blood in an artery is modeled by the nonlinear equation $Q=K\left(P-\dfrac{Q^{2}}{R}\right)$ where Q is the flow rate (mL / s) P= 100 mmHg is the blood pressure, k = 0.05 and R = 10. Determine the possible real values of the flow rate Q accurate to two decimal places.}
\q{509}{In a nonlinear LC circuit, the equilibrium charge q on the capacitor satisfies q$^{3}$-4q+2=0. Find the equilibrium values of charge stored on the capacitor lying in the interval [0, 1] using Regula-Falsi method accurate to three decimal places.}
\q{510}{In a certain cryptographic system, the decoding speed $v$ (in units per second) is related to the security delay function where $d(v)=\dfrac{v^{2}}{20+v}$ , where d(v) represents the decoding delay in milliseconds. For secure transmission, the systém requires a decoding delay of 50 milliseconds. Use Newton-Raphson Method with initial guess $v_{0}=35$ to determine the decoding speed v that satisfies d(v) = 50 , correct to three decimal places.}
\q{511}{In a secure communication system, the phase shift $\theta$ (in radians) of an encrypted signal is modeled by the nonlinear function $f(\theta)=\sin\theta-0.5$ For proper synchronization of the encrypted signal, the phase shift must satisfy $f(\theta)=0$. Use the Newton-Raphson Method with initial guess $\theta_{0}=0.5$ to determine the required phase shift $\theta$, correct to three decimal places.}
\q{512}{In a nonlinear electric circuit, the voltage x (in volts) across a component is modeled by $f(x)=-0.1x^{2}+2x+1$ . Use the Newton-Raphson Method with initial guess $x_{0}=15$ to estimate the operating voltage at which the circuit reaches equilibrium. The equilibrium voltage is the voltage at which the net change in the circuit becomes zero, i.e., the circuit reaches a stable operating point.}
\q{513}{In a chemical reaction, the concentration x (in moles per liter) of a product at equilibrium satisfies the the relation $xe^{x}=1$ . Find the equilibrium concentration correct to two decimal places using Newton-Raphson Method.}
\q{514}{The electrical activity of the heart is represented by a variable x, which measures the signal strength of the heartbeat. At steady-state, the signal satisfies the nonlinear equation: $x^{3}-x-1=0$}
\q{515}{In a circuit, the voltage $V$ across a component satisfies the equation: $\sin V-0.5V=0$ . Find a non-zero real value of $V$ correct to two decimal places using Regula-Falsi method.}
\q{516}{In a cryptosystem, a key parameter x satisfies the equation: $x^{5}-10x+5=0$ . Find a real value of x correct to two decimal places using the Regula-Falsi method.}
\q{517}{Find the principal value of the following: $\tan^{-1}\!\left[3\cos\left\{3\cos^{-1}\!\left(\dfrac{\sqrt{3}}{2}\right)\right\}\right]$}
\q{518}{Find the principal value of the following: $\csc^{-1}\!\left[\dfrac{4}{3}\cos\left\{\tan^{-1}\!\left(\dfrac{1}{\sqrt{3}}\right)\right\}\right]$}
\q{519}{Show that: $\cos^{-1}\dfrac{8}{17}=\sin^{-1}\dfrac{15}{17}$}
\q{520}{Show that: $\csc^{-1}\dfrac{5}{4}=\cos^{-1}\dfrac{3}{5}$}
\q{521}{Show that: $2\tan^{-1}\dfrac{3}{4}=\tan^{-1}\dfrac{24}{7}$}
\q{522}{Show that: $2\sin^{-1}\dfrac{7}{25}=\sin^{-1}\dfrac{336}{625}$}
\q{523}{Without using calculator find the value of the following: $\cot\!\left(\dfrac{1}{2}\sec^{-1}\dfrac{2\sqrt{6}}{3}\right)$}
\q{524}{Without using calculator find the value of the following: $\csc\!\left(\dfrac{1}{2}\cot^{-1}\dfrac{4}{3}\right)$}
\q{525}{Find value of: $\csc\left[\sin^{-1}\!\left(\dfrac{\sqrt{3}}{2}\right)\right]$}
\q{526}{Prove that: $\sin^{-1}x+\cos^{-1}x=\dfrac{\pi}{2},x\in[-1,1]$}
\q{527}{Prove that: $\csc^{-1}x+\sec^{-1}x=\dfrac{\pi}{2},x\in(-\infty,-1]\cup[1,\infty)$}
\q{528}{Prove that: $\tan^{-1}x+\cot^{-1}x=\dfrac{\pi}{2},x\in(-\infty,\infty)$}
\q{529}{Show that $\cos^{-1}\dfrac{7}{25}=\csc^{-1}\dfrac{25}{24}$}
\q{530}{Without using calculator find the value of $\sin\left[\dfrac{1}{2}\tan^{-1}\dfrac{4}{3}\right]$}
\q{531}{Prove the following: $\sin^{-1}\dfrac{7}{25}+\sin^{-1}\dfrac{15}{17}=\cos^{-1}\dfrac{87}{425}$}
\q{532}{Prove the following: $\cos^{-1}\dfrac{7}{25}-\cos^{-1}\dfrac{5}{13}=\cos^{-1}\dfrac{323}{325}$}
\q{533}{Prove the following: $\tan^{-1}\dfrac{3}{4}+\tan^{-1}\dfrac{3}{5}-\tan^{-1}\dfrac{8}{19}=\dfrac{\pi}{4}$}
\q{534}{Prove the following: $\sin^{-1}\dfrac{12}{13}-\sin^{-1}\dfrac{3}{5}=\sin^{-1}\dfrac{33}{65}$}
\q{535}{Prove the following: $\sin^{-1}\dfrac{4}{5}+\cos^{-1}\dfrac{12}{13}=\sin^{-1}\dfrac{63}{65}$}
\q{536}{Prove the following: $\sin^{-1}\dfrac{3}{5}+2\tan^{-1}\dfrac{1}{4}=\cos^{-1}\dfrac{36}{85}$}
\q{537}{Prove the following: $\sin^{-1}\dfrac{1}{\sqrt{2}}+\sin^{-1}\dfrac{1}{\sqrt{5}}+\sin^{-1}\dfrac{1}{\sqrt{10}}=\dfrac{\pi}{2}$}
\q{538}{Prove the following: $\tan^{-1}\dfrac{1}{7}+2\tan^{-1}\dfrac{1}{3}=\dfrac{\pi}{4}$}
\q{539}{Prove the following: $\tan^{-1}\dfrac{10}{3}-\tan^{-1}\dfrac{7}{13}+\tan^{-1}\dfrac{9}{5}-\tan^{-1}\dfrac{2}{7}=\dfrac{\pi}{2}$}
\q{540}{Prove the following: $\sec^{2}\!\left(\tan^{-1}\dfrac{1}{4}\right)+\csc^{2}\!\left(\cot^{-1}\dfrac{2}{5}\right)=\dfrac{889}{400}$}
\q{541}{Prove the following: $\tan^{2}(\sec^{-1}3)+\cot^{2}(\csc^{-1}4)=23$}
\q{542}{Prove the following: $2\tan^{-1}x=\sin^{-1}\!\left(\dfrac{2x}{1+x^{2}}\right),-1\leq x\leq1$}
\q{543}{Prove the following: $\tan(\sin^{-1}x)=\dfrac{x}{\sqrt{1-x^{2}}},-1<x<1$}
\q{544}{Prove the following: $\tan^{-1}x+\cot^{-1}(x+1)=\tan^{-1}(x^{2}+x+1)$}
\q{545}{Prove the following: $2\tan^{-1}\!\left(\tan\dfrac{\theta}{2}\tan\dfrac{\phi}{2}\right)=\cos^{-1}\!\left(\dfrac{\cos\theta+\cos\phi}{1+\cos\theta\cos\phi}\right)$}
\q{546}{If $\tan^{-1}x+\tan^{-1}y+\tan^{-1}z=\pi$, then show that x+y+z=xyz.}
\q{547}{Find the value of x if, $\tan^{-1}\!\left(\dfrac{2x-1}{3x-1}\right)+\tan^{-1}\!\left(\dfrac{2x+1}{3x+1}\right)=\dfrac{\pi}{4}$}
\q{548}{If $\sin^{-1}\dfrac{2p}{1+p^{2}}+\cos^{-1}\dfrac{1-q^{2}}{1+q^{2}}=\tan^{-1}\dfrac{2y}{1-y^{2}}$, show that y=}
\q{549}{Prove that $\sin^{-1}\!\left(\dfrac{5}{13}\right)+\sin^{-1}\!\left(\dfrac{12}{13}\right)=\dfrac{\pi}{2}$}
\q{550}{Prove that $\tan^{-1}\!\left(\dfrac{x}{y}\right)-\tan^{-1}\!\left(\dfrac{x-y}{x+y}\right)=\dfrac{\pi}{4}$}
\q{551}{Solve the equation $\tan^{-1}\!\left(\dfrac{x-1}{x+1}\right)+\tan^{-1}\!\left(\dfrac{2x}{x+1}\right)=\dfrac{\pi}{4}$}
\q{552}{Solve the following trigonometric eqautions: $\tan^{2}x-(\sqrt{3}+1)\tan x=-\sqrt{3}$}
\q{553}{Solve the following trigonometric eqautions: $\sqrt{2\sin^{2}x+1}=\sin x+\cos x$}
\q{554}{Solve the following trigonometric eqautions: $\sin^{3}x\cos x+\cos^{3}x\sin x=\dfrac{1}{2}$}
\q{555}{Solve the following trigonometric eqautions: $\sin^{3}2x\cos 2x-\cos^{3}2x\sin 2x=\dfrac{\sqrt{2}}{8}$}
\q{556}{Solve the following trigonometric eqautions: $\sin 3x\cos x+\cos 3x\sin x=\dfrac{1}{\sqrt{2}}$}
\q{557}{Solve the following trigonometric eqautions: $\dfrac{\sin x}{1+\cos x}+\dfrac{1+\cos x}{\sin x}=4$}
\q{558}{Solve the following trigonometric eqautions: $\dfrac{\tan\left(7x+\frac{\pi}{6}\right)+\tan 3x}{1-\tan\left(7x+\frac{\pi}{6}\right)\tan 3x}=1$}
\q{559}{Solve the following trigonometric eqautions: $\dfrac{\cot\left(2x+\frac{\pi}{12}\right)-\cot\left(5x+\frac{\pi}{12}\right)}{1+\cot\left(2x+\frac{\pi}{12}\right)\cot\left(5x+\frac{\pi}{12}\right)}=\dfrac{1}{\sqrt{3}}$}
\q{560}{Solve the following trigonometric equations: $\dfrac{\sin^{3}x+\cos^{3}x}{\sin x+\cos x}-\dfrac{\sin x}{\sqrt{1+\tan^{2}x}}=\dfrac{1}{2}$}
\q{561}{Solve the following trigonometric eqautions: $\sin 3x-\sin 5x+\sin 7x=0$}
\q{562}{Solve the following trigonometric eqautions: $\cos 2x-\cos 6x=\sin 4x$}
\q{563}{Solve the following trigonometric equations: sin x+ sin 3x + sin5x+sin 7x=0}
\q{564}{Solve $\sin x=-\dfrac{\sqrt{3}}{2}$}
\q{565}{Solve $2\sec^{2}x-\dfrac{8}{3}=0$}
\q{566}{Solve $\sqrt{3}$cotx-cscx-1=0}
\q{567}{Solve graphically: $3\tan x=2x+1,\ x\in\left(-\dfrac{\pi}{2},\dfrac{\pi}{2}\right)$}
\q{568}{Solve graphically: $\sec(x+3)=-2x+\dfrac{1}{3},\ x\in[-\pi,\pi]$}
\q{569}{Solve graphically: $\sin 3x=\dfrac{x}{2}$}
\q{570}{A hiker pulls a sled using a rope that exerts a force of 80 newtons over a distance of 25 meters. The total work done is 1200 joules. At what angle to the direction of motion is the force being applied?}
\q{571}{A radio tower stands on flat ground with a small antenna mounted at its top. From a point 50 meters away from the base of the tower. The antenna subtends the same angle as a 1 meter vertical rod placed at the base of the tower. The total height of the tower and antenna is 10 meters.\par (a) Find the height of the tower (without the antenna).\par (b) Calculate the subtended angle.}
\q{572}{A person standing 20 meters away from the base of a tree measures the angle of elevation to the top of the tree as 30°.\par (a) Find the height of the tree.\par (b) If the person walks 10 meters closer to the tree, what will be the new angle of elevation to the top of the tree?}
\q{573}{Solve the equation $x=\dfrac{x}{2}$ graphically over the interval $[-\pi,\pi]$.}
\q{574}{Solve the equation $\cos x=\dfrac{x}{8}$ graphically.}
\q{575}{Solve the equation $x=\dfrac{x}{2}$ graphically in $\left(-\dfrac{3\pi}{2},\dfrac{3\pi}{2}\right)$}
\q{576}{A lighthouse with a 15 meter tall lantern stands atop a cliff. From a ship 120 meters offshore, the lantern subtends the same angle as a 1.5 meter buoy placed at the base of the cliff.\par (a) Find the height of the cliff.\par (b) Calculate the subtended angle $\theta$}
\q{577}{A projectile is launched from the ground with an initial velocity $v_{\circ}$ 30m/s. The range $R$ of the projectile is given by the formula $R=\dfrac{v_{\circ}^{2}}{g}\sin 2\theta$, where $\theta$ is the launch angle to the horizontal $g=9.8\ m/s^{2}$ is the acceleration due to gravity.\par (a) Find the angle $\theta$ (in degrees) required for the projectile to hit a target exactly 80 meters away.\par (b) If the maximum range achievable is $R_{max}=\dfrac{v_{\circ}^{2}}{g}$ determine the angle $\theta$ that gives the maximum range. Justify your answer mathematically.}
\q{578}{Given $x(t)=\dfrac{1}{t+4},\ y(t)=\ln\sqrt{t},\ z(t)=\sec^{-1}(t)$, construct r(t).}
\q{579}{Find the Find the domain and range of the following vector-valued functions: $r(t)=\sin(t)\,\hat{i}+\cos(t)\,\hat{j}+t\hat{k}$}
\q{580}{Determine which of the following functions are scalar functions or vector-valued functions. Give appropriate reasoning to support the answer: $f(t)=(3t^{2}-5)^{2}$}
\q{581}{Determine which of the following functions are scalar functions or vector-valued functions. Give appropriate reasoning to support the answer: $h(t)=\langle t^{2},\,e^{t},\,\ln\sqrt{t}\rangle$}
\q{582}{Determine which of the following functions are scalar functions or vector-valued functions. Give appropriate reasoning to support the answer: $g(t)=\tan(t)+\csc(t)$}
\q{583}{Determine which of the following functions are scalar functions or vector-valued functions. Give appropriate reasoning to support the answer: $j(t)=\left\langle\dfrac{1}{t^{2}+1},\dfrac{1}{t^{2}+1}\right\rangle$}
\q{584}{Find the domain and range of the vector-valued function $r(t)=\sqrt{4-t^{2}}\,\hat{i}+\dfrac{1}{t-1}\hat{j}$.}
\q{585}{Find the first-order derivatives of the following vector-valued functions with respect to t. $\boldsymbol{r}(t)=\dfrac{2}{\sqrt{1-t^{2}}}\hat{i}+\dfrac{\sqrt{3}}{1+t^{2}}\hat{k}$}
\q{586}{Find the first-order derivatives of the following vector-valued functions with respect to t at the specified values of t. $r(t)=\left(\dfrac{t^{2}+1}{\sin t}\right)\hat{i}+(\sqrt{t}\,e^{t})\hat{j}$ at $t=\dfrac{\pi}{2}$}
\q{587}{Find the first-order derivatives of the following vector-valued functions with respect to t at the specified values of t. $\boldsymbol{r}(t)=\left(t\sqrt{t+2}\right)\hat{i}+\ln\!\left(t+\sqrt{t^{2}+1}\right)\hat{j}$ at $t=0$}
\q{588}{Find the first-order derivatives of the following vector-valued functions with respect to t at the specified values of t. $\boldsymbol{r}(t)=\arctan(t)\,\hat{i}+\dfrac{\sin(t^{2})}{t}\hat{j}\ \sqrt[3]{1+t}\ \hat{k}$ at $t=0$}
\q{589}{Find the first-order derivatives of the following vector-valued functions with respect to t at the specified values of t. $\boldsymbol{r}(t)=te^{-t}\,\hat{i}+\cos(3t)\,\hat{j}+\ln(1+t^{2})\,\hat{k}$ at $t=0$}
\q{590}{Find the first derivative of $r(t)=\sec t\hat{i}+\tan t\hat{j}-\cos t\hat{k}$ at $t=\dfrac{\pi}{4}$}
\q{591}{Find the velocity, acceleration, speed, and direction of motion at a specified value of 1. Interpret the results by describing the particle's motion at that instant, including how fast and in what direction it is moving, whether it is speeding up or slowing down, and how its path is curving.\par $\boldsymbol{r}(t)=2\sin t\,\hat{i}+3t\hat{j}$ at $t=\dfrac{\pi}{2}$}
\q{592}{Find the velocity, acceleration, speed, and direction of motion at a specified value of 1. Interpret the results by describing the particle's motion at that instant, including how fast and in what direction it is moving, whether it is speeding up or slowing down, and how its path is curving.\par $\boldsymbol{r}(t)=e^{t}\hat{i}+e^{-t}\hat{j}$ at $t=0$}
\q{593}{Find the velocity, acceleration, speed, and direction of motion at a specified value of 1. Interpret the results by describing the particle's motion at that instant, including how fast and in what direction it is moving, whether it is speeding up or slowing down, and how its path is curving.\par $\boldsymbol{r}(t)=\cos t\hat{i}+\sin t\hat{j}+2t\hat{k}$ at $t=\pi$}
\q{594}{Find the velocity, acceleration, speed, and direction of motion at a specified value of 1. Interpret the results by describing the particle's motion at that instant, including how fast and in what direction it is moving, whether it is speeding up or slowing down, and how its path is curving.\par $\boldsymbol{r}(t)=t^{2}\hat{i}+t\hat{j}+t^{3}\hat{k}$ at $t=2$}
\q{595}{Find the velocity, acceleration, speed, and direction of motion at a specified value of 1. Interpret the results by describing the particle's motion at that instant, including how fast and in what direction it is moving, whether it is speeding up or slowing down, and how its path is curving.\par $r(t)=e^{t}\sin t\hat{i}+(t^{3}-2t)\hat{j}+\arctan(2t)\hat{k}$ at $t=0$}
\q{596}{A cannonball is fired from the ground at an angle $30^{o}$ with an initial velocity of 100 m/s. Assuming gravity acts downward $(g=9.8m/(s^{2}))$ and neglecting air resistance,\par (a)Formulate the position vector r(t) of the cannonball as function of time.\par (b) Find the velocity vector at t = 5 sec.\par (c) Determine the time of the flight (when the cannonball hits the ground).}
\q{597}{A comet orbits the Sun in an elliptical path such that its position (in AU) at time \textbf{t} (in years) is described by, x = 2cos t, y==sin \textbf{t}.\par (a) Write the position vector r(t)\par (b) Find the velocity and speed at $t=\dfrac{\pi}{2}$ .\par (c) Show that the acceleration vector always points toward the origin.}
\q{598}{A robotic arm moves such that: Its x-position increases linearly at 2m /s Its y-position follows $y(t)=t^{2}$, Its z-position is given by $z(t)=1-e^{-t}$.\par (a) Construct the position vector T(t)\par (b) Find the velocity vector at t = 1 sec\par (c) Compute the acceleration at t = 1 sec}
\q{599}{A particle moves along a circular path, described by the position vector $\boldsymbol{r}(t)=R\cos\omega t\hat{i}+R\sin\omega t\hat{j}$ . Prove that its velocity is perpendicular to its position vector.}
\q{600}{A Ferris wheel of radius 20 m rotates with an angular speed of 0.2 rad/sec.\par (a) Construct its position vector r(t)\par (b) Find the passenger's velocity at t =10sec\par (c) Show that the acceleration is centripetal (toward the center)}
\q{601}{A quarterback throws a football with a position vector $r(t)=20t\hat{i}+(20t-4.9t^{2})\hat{j}$ .\par (a) Find the time of flight before it hits the ground.\par (b) What is the maximum height reached?\par (c) At which time did football reach its maximum height?}
\q{602}{A small metal part moves on a conveyor belt in a factory. The belt transports the part eastward (x-direction) while simultaneously vibrating sideways (north-south, y-direction) to help sort items into bins. The position of the part at time t seconds is given by $\boldsymbol{r}(t)=3t^{2}\,\hat{i}+2\sin t\hat{j}$ sint in meters. At a specific moment, $t=\dfrac{\pi}{2}$ an engineer needs to analyze the part's motion to ensure it enters the correct bin. Hence, find the following quantities at $t=\dfrac{\pi}{2}$: (a) velocity (b) acceleration (c) speed (d) direction of motion.}
\q{603}{A drone is programmed to fly in a helical (spiral) path around a large cylindrical water tower to inspect its outer surface for cracks or rust. The drone maintains a constant horizontal distance of 2 meters from the tower's center while climbing upward at a steady rate. The drone's position in 3D space at time t seconds is given by:\par $\boldsymbol{r}(t)=2\cos t\hat{i}+2\sin t\hat{j}+3t\hat{k}$\par Here:\par x(t) = 2cos t east-west position (amplitude 2 m)\par y(t) = 2sin t north-south position (amplitude 2 m)\par z(t) = 3t Vertical height (climbs 3m / s )\par At a specific moment $t=\pi$ seconds, the drone operator needs to analyze its motion to ensure smooth inspection. Find the following quantities at $t=\pi$ (a) velocity, (b) acceleration, (c) speed, and (d) direction of motion at $t=\pi$ .}
\end{document}
